\documentclass[a4paper,12pt,oneside]{report}
% ======================================================================
%  Định dạng báo cáo — theo mẫu Tiểu luận 03 (Times 13pt, giãn dòng 1.5)
% ======================================================================
\usepackage[a4paper,left=3cm,right=2cm,top=2.5cm,bottom=2.5cm,headheight=14pt]{geometry}
\usepackage{fontspec}
\usepackage{unicode-math}
\setmainfont{TeX Gyre Termes}[Scale=1.0833]          % 12pt × 1,0833 ≈ 13pt
\setmathfont{TeX Gyre Termes Math}[Scale=1.0833]
\setmonofont{SourceCodePro}[Extension=.otf,UprightFont=*-Regular,BoldFont=*-Bold,ItalicFont=*-RegularIt,BoldItalicFont=*-BoldIt,Scale=0.82]
\usepackage{polyglossia}
\setdefaultlanguage{vietnamese}
\setotherlanguage{english}  % để biblatex dùng chuỗi tiếng Anh (đã Việt hóa) trong danh mục tài liệu

\usepackage{setspace}
\setstretch{1.4}
\setlength{\parindent}{1cm}
\setlength{\parskip}{3pt}
\usepackage{indentfirst}
\usepackage{microtype}
\usepackage{ragged2e}
\setlength{\emergencystretch}{3em}

\usepackage{amsmath,amsthm}
\usepackage{graphicx}
\graphicspath{{figures/}}
\usepackage{float}
\usepackage{booktabs,longtable,tabularx,multirow,array,makecell}
\renewcommand{\arraystretch}{1.18}
\usepackage{xcolor}
\definecolor{codebg}{HTML}{F4F5F7}
\definecolor{accent}{HTML}{1D4ED8}
\usepackage[hidelinks]{hyperref}
\usepackage{xurl}
\usepackage[htt]{hyphenat}  % cho phép ngắt dòng trong \texttt (tên hàm, đường dẫn dài)
\usepackage{adjustbox}
\usepackage{enumitem}
\setlist{itemsep=1pt,topsep=3pt,parsep=0pt,leftmargin=1.2cm}
\usepackage{tikz}
\usetikzlibrary{positioning,arrows.meta,shapes.geometric,shapes.symbols,calc,fit,backgrounds}

% ---------------------------------------------------------------- tiêu đề chương, mục
\usepackage{titlesec}
\titleformat{\chapter}[block]{\centering\bfseries\fontsize{15}{19}\selectfont}
  {\MakeUppercase{\chaptertitlename}\ \thechapter.}{0.5em}{\MakeUppercase}
\titlespacing{\chapter}{0pt}{-10pt}{18pt}
\titleformat{\section}{\bfseries\fontsize{13.5}{17}\selectfont}{\thesection.}{0.5em}{}
\titleformat{\subsection}{\bfseries\itshape\fontsize{13}{16}\selectfont}{\thesubsection.}{0.5em}{}
\titleformat{\subsubsection}{\itshape\fontsize{13}{16}\selectfont}{\thesubsubsection.}{0.5em}{}
\titlespacing{\section}{0pt}{10pt}{4pt}
\titlespacing{\subsection}{0pt}{8pt}{3pt}
\titlespacing{\subsubsection}{0pt}{6pt}{2pt}
\setcounter{secnumdepth}{3}
\setcounter{tocdepth}{2}
\renewcommand{\chaptertitlename}{Chương}

% chương không đánh số (Mở đầu, Kết luận, TLTK) vẫn vào mục lục
\newcommand{\unnumberedchapter}[1]{%
  \chapter*{#1}\addcontentsline{toc}{chapter}{#1}\markboth{#1}{#1}}

% ---------------------------------------------------------------- mục lục
\usepackage{tocloft}
\renewcommand{\contentsname}{\centering MỤC LỤC}
\renewcommand{\listfigurename}{\centering DANH MỤC HÌNH VẼ}
\renewcommand{\listtablename}{\centering DANH MỤC BẢNG BIỂU}
\renewcommand{\cftchapfont}{\bfseries}
\renewcommand{\cftchappagefont}{\bfseries}
\renewcommand{\cftchappresnum}{CHƯƠNG }
\setlength{\cftchapnumwidth}{6.6em}
\renewcommand{\cftchapaftersnum}{.}
\renewcommand{\cftsecaftersnum}{.}
\renewcommand{\cftsubsecaftersnum}{.}
\renewcommand{\cftchapleader}{\cftdotfill{\cftdotsep}}
\renewcommand{\cftfigpresnum}{Hình }
\renewcommand{\cfttabpresnum}{Bảng }
\setlength{\cftfignumwidth}{5.2em}
\setlength{\cfttabnumwidth}{5.2em}
\renewcommand{\cftfigaftersnum}{.}
\renewcommand{\cfttabaftersnum}{.}

% ---------------------------------------------------------------- chú thích hình, bảng
\usepackage{caption}
\DeclareCaptionLabelSeparator{dotspace}{.\ }
\captionsetup{labelsep=dotspace,font={small},justification=centering,singlelinecheck=false}
\captionsetup[figure]{name=Hình,font={small,it},skip=4pt}
\captionsetup[table]{name=Bảng,font={small,bf},position=top,skip=4pt}
\renewcommand{\thefigure}{\thechapter.\arabic{figure}}
\renewcommand{\thetable}{\thechapter.\arabic{table}}

% ---------------------------------------------------------------- đoạn mã
\usepackage[newfloat]{minted}
\setminted{fontsize=\footnotesize,breaklines,bgcolor=codebg,frame=none,tabsize=4,
           baselinestretch=1.05,style=friendly,autogobble}
\SetupFloatingEnvironment{listing}{name=Đoạn mã,placement=H}
\captionsetup[listing]{font={small,bf},position=top,justification=raggedright,skip=2pt}
\renewcommand{\thelisting}{\thechapter.\arabic{listing}}
\newenvironment{code}[1]{\begin{listing}[H]\caption{#1}}{\end{listing}}

% ---------------------------------------------------------------- trang
\usepackage{fancyhdr}
\pagestyle{fancy}
\fancyhf{}
\fancyfoot[C]{\thepage}
\renewcommand{\headrulewidth}{0pt}
\fancypagestyle{plain}{\fancyhf{}\fancyfoot[C]{\thepage}\renewcommand{\headrulewidth}{0pt}}

% ---------------------------------------------------------------- tài liệu tham khảo
\usepackage[backend=biber,style=ieee,sorting=none,maxnames=4,minnames=1,language=english,autolang=none]{biblatex}
\addbibresource{refs.bib}
\DefineBibliographyStrings{english}{andothers={và cộng sự},and={và},in={trong},edition={ấn bản},volume={tập},jourvol={tập},number={số},pages={tr.},page={tr.}}

% ---------------------------------------------------------------- tiện ích
\newcommand{\eq}[1]{(\ref{#1})}
\newcommand{\figref}[1]{Hình~\ref{#1}}
\newcommand{\tabref}[1]{Bảng~\ref{#1}}
\newcommand{\coderef}[1]{Đoạn mã~\ref{#1}}
\newcommand{\weburl}[1]{{\small\url{#1}}}
\newcommand{\api}[1]{\texttt{\small #1}}
\newcommand{\figw}[3][\linewidth]{%
  \begin{figure}[H]\centering\includegraphics[width=#1]{#2}\caption{#3}\label{fig:#2}\end{figure}}
\newcolumntype{L}[1]{>{\raggedright\arraybackslash}p{#1}}
\newcolumntype{C}[1]{>{\centering\arraybackslash}p{#1}}
\newcolumntype{Y}{>{\raggedright\arraybackslash}X}
\newtheorem*{dinhnghia}{Định nghĩa}

% tên miền triển khai (Cloudflare Tunnel) — đổi ở đây là đổi toàn bộ báo cáo
\newcommand{\deploydomain}{agent.nits.io.vn}
\newcommand{\agent}[1]{\texttt{#1}}


\begin{document}

% ====================================================================== BÌA (giữ nguyên mẫu bìa của sinh viên)
\pagenumbering{Alph}
\begin{titlepage}
\begin{tikzpicture}[remember picture,overlay]
  \draw[line width=1.6pt] ($(current page.north west)+(2.3cm,-1.8cm)$) rectangle ($(current page.south east)+(-1.5cm,1.8cm)$);
  \draw[line width=0.5pt] ($(current page.north west)+(2.45cm,-1.95cm)$) rectangle ($(current page.south east)+(-1.65cm,1.95cm)$);
\end{tikzpicture}
\centering
\setstretch{1.25}
\vspace*{0.4cm}
{\bfseries HỌC VIỆN CÔNG NGHỆ BƯU CHÍNH VIỄN THÔNG\\[2pt] KHOA CÔNG NGHỆ THÔNG TIN 1\\[-4pt]}
\rule{5.2cm}{0.5pt}

\vspace{1.0cm}
\includegraphics[height=3.4cm]{ptit_logo}

\vspace{1.0cm}
{\fontsize{19}{25}\selectfont\bfseries PROJECT 01\\[4pt] PHÁT TRIỂN HỆ THỐNG THÔNG MINH\par}
\vspace{1.1cm}
{\fontsize{14.5}{20}\selectfont\bfseries HỆ ĐA AGENT CHO BỆNH TIỂU ĐƯỜNG\par}
\vspace{0.4cm}
{\itshape Deep Learning, Knowledge Base (Neo4j) và RAG,\\ Reinforcement Learning kết hợp luật, hội đồng tranh biện\par}

\vspace{1.6cm}
\begin{tabular}{@{}l@{\hspace{1.2cm}}l@{}}
\textbf{Giảng viên hướng dẫn:} & PGS. TS. Trần Đình Quế \\
\textbf{Sinh viên thực hiện:} & Nguyễn Tự Kiên \\
\textbf{Mã sinh viên:} & B23DCCN465 \\
\textbf{Nhóm lớp:} & CT01 \\
\textbf{Nhóm:} & 07 \\
\end{tabular}

\vfill
{\bfseries Hà Nội, tháng 11 năm 2026\par}
\vspace{0.6cm}
\end{titlepage}

% ====================================================================== PHẦN ĐẦU
\pagenumbering{roman}
\setcounter{page}{2}

\unnumberedchapter{TÓM TẮT BÁO CÁO}
Báo cáo này trình bày kiến trúc tổng thể và kết quả triển khai thực nghiệm của \textbf{Hệ thống Đa Tác tử Lai Hỗ trợ Quyết định Lâm sàng Bệnh Đái tháo đường} (\emph{Hybrid Multiagent System Architecture for Diabetes Decision Support}). Hệ thống được xây dựng và mở rộng trên nền tảng hai thành phần cốt lõi đã phát triển: (i) tổ hợp các mô hình Học sâu (\emph{Deep Learning}) và Học máy (\emph{Machine Learning}) dự báo nguy cơ đái tháo đường và biến chứng được huấn luyện trên dữ liệu thực tế quy mô lớn (CDC NHANES 2011--2023, CDC BRFSS 2021--2023 và bộ dữ liệu nhân trắc người Việt), và (ii) hệ thống Đồ thị Tri thức (\emph{Knowledge Graph}) trên Neo4j kết hợp Sinh văn bản tăng cường truy xuất lai (\emph{Hybrid Retrieval-Augmented Generation -- RAG}) trên Qdrant và SQLite FTS5 được chuẩn hóa từ bốn nguồn tri thức y khoa: Hiệp hội Đái tháo đường Hoa Kỳ (ADA 2024), Quyết định 5481/QĐ-BYT của Bộ Y tế, Hệ thống Y tế Vinmec và Dược điển Nhà thuốc FPT Long Châu.

Kiến trúc đề xuất bao gồm các tác tử chuyên biệt (\emph{specialized agents}) đảm nhiệm việc kiểm định và chuẩn hóa dữ liệu bệnh nhân (\emph{Patient Data Agent}), suy diễn dự báo bằng Học sâu và Học máy (\emph{DL Prediction Agent}), suy luận dựa trên luật lâm sàng và phản biện đồng thuận (\emph{Rule-Based \& Consensus Critic Agent}), truy xuất tri thức đồ thị và văn bản (\emph{KG/RAG Agent}), giải thích lâm sàng có dẫn chứng bằng mô hình ngôn ngữ lớn cục bộ Ollama/Qwen2.5 (\emph{Explanation Agent}), giám sát và kiểm toán vết thực thi (\emph{Monitoring Agent}), dưới sự điều phối tập trung của tác tử điều phối (\emph{Coordinator Agent}). Ngoài ra, tác tử Học tăng cường (\emph{Reinforcement Learning Agent}) sử dụng thuật toán Q-Learning kết hợp Cổng an toàn lâm sàng (\emph{Clinical Safety Gate}) được tích hợp như một mô-đun nâng cao hoạt động trong môi trường mô phỏng nhằm gợi ý lộ trình can thiệp bậc thang an toàn.

Tài liệu này định nghĩa chi tiết kiến trúc phân tầng sáu lớp, phân loại trách nhiệm tác tử, mô hình hóa toán học, luồng tương tác tuần tự và song song, hợp đồng giao tiếp dữ liệu, cơ chế tích hợp mô hình -- tri thức, lộ trình triển khai chín giai đoạn cùng các nguyên tắc an toàn y khoa. Hệ thống là một nguyên mẫu nghiên cứu và giáo dục (\emph{educational and research prototype}), đóng vai trò công cụ hỗ trợ sàng lọc và tham khảo lâm sàng có giải trình, không thay thế quyết định chẩn đoán hay kê đơn trực tiếp của bác sĩ chuyên khoa.

\clearpage
\begin{spacing}{1.15}
\tableofcontents
\end{spacing}
\clearpage
\chapter*{DANH MỤC TỪ VIẾT TẮT}
\addcontentsline{toc}{chapter}{DANH MỤC TỪ VIẾT TẮT}
\begin{spacing}{1.1}
\small
\begin{longtable}{@{}L{2.3cm}L{6.3cm}L{6.0cm}@{}}
\toprule
\textbf{Viết tắt} & \textbf{Tiếng Anh} & \textbf{Tiếng Việt} \\
\midrule
\endhead
ADA & American Diabetes Association & Hiệp hội Đái tháo đường Hoa Kỳ \\
AUC & Area Under the (ROC) Curve & Diện tích dưới đường ROC \\
BM25 & Best Matching 25 & Hàm xếp hạng BM25 \\
BMI & Body Mass Index & Chỉ số khối cơ thể \\
BRFSS & Behavioral Risk Factor Surveillance System & Hệ thống Giám sát Yếu tố Nguy cơ Hành vi (CDC) \\
CGM & Continuous Glucose Monitoring & Theo dõi đường huyết liên tục \\
CNN & Convolutional Neural Network & Mạng nơ-ron tích chập \\
DL & Deep Learning & Học sâu \\
ĐTĐ & -- & Đái tháo đường \\
eGFR & Estimated Glomerular Filtration Rate & Độ lọc cầu thận ước tính \\
FIPA-ACL & FIPA Agent Communication Language & Ngôn ngữ giao tiếp tác tử FIPA \\
FPG & Fasting Plasma Glucose & Glucose huyết tương lúc đói \\
FTS & Full-Text Search & Tìm kiếm toàn văn \\
GBDT & Gradient Boosted Decision Trees & Cây quyết định tăng cường gradient \\
GLP-1 RA & Glucagon-Like Peptide-1 Receptor Agonist & Đồng vận thụ thể GLP-1 \\
HbA1c & Glycated Hemoglobin & Huyết sắc tố gắn glucose \\
HNSW & Hierarchical Navigable Small World & Đồ thị thế giới nhỏ phân cấp \\
KB & Knowledge Base & Cơ sở tri thức \\
KG & Knowledge Graph & Đồ thị tri thức \\
LLM & Large Language Model & Mô hình ngôn ngữ lớn \\
LSTM & Long Short-Term Memory & Mạng bộ nhớ ngắn-dài hạn \\
MAS & Multi-Agent System & Hệ đa tác tử (hệ đa agent) \\
MDP & Markov Decision Process & Quá trình quyết định Markov \\
ML & Machine Learning & Học máy \\
MLP & Multilayer Perceptron & Mạng nơ-ron nhiều lớp \\
NHANES & National Health and Nutrition Examination Survey & Khảo sát Sức khỏe và Dinh dưỡng Quốc gia Hoa Kỳ \\
OGTT & Oral Glucose Tolerance Test & Nghiệm pháp dung nạp glucose đường uống \\
QWK & Quadratic Weighted Kappa & Hệ số đồng thuận Kappa trọng số bậc hai \\
RAG & Retrieval-Augmented Generation & Sinh văn bản tăng cường truy xuất \\
RL & Reinforcement Learning & Học tăng cường \\
RRF & Reciprocal Rank Fusion & Gộp xếp hạng nghịch đảo \\
SGLT2i & Sodium-Glucose Cotransporter-2 Inhibitor & Thuốc ức chế kênh đồng vận chuyển SGLT2 \\
\bottomrule
\end{longtable}
\end{spacing}

\clearpage
\begin{spacing}{1.1}
\listoffigures
\clearpage
\listoftables
\end{spacing}
\clearpage

% ====================================================================== NỘI DUNG (11 chương theo đúng form của giảng viên)
\pagenumbering{arabic}
\chapter{Giới thiệu}
\label{chap:intro}

\section{Bối cảnh dự án (Project Background)}
\label{sec:project-background}

Đái tháo đường típ~2 (\emph{Type~2 Diabetes Mellitus} -- ĐTĐ típ~2) là một trong những bệnh lý chuyển hóa mạn tính có tốc độ gia tăng nhanh nhất trên toàn cầu cũng như tại Việt Nam~\cite{idf2021,byt5481}. Đặc thù lâm sàng của bệnh là diễn tiến âm thầm trong nhiều năm ở giai đoạn tiền đái tháo đường (\emph{prediabetes}) trước khi xuất hiện các biến chứng mạch máu nhỏ (võng mạc, thận, thần kinh ngoại biên) và mạch máu lớn (nhồi máu cơ tim, đột quỵ)~\cite{ada2024}. Trong thực tế hỗ trợ quyết định lâm sàng (\emph{Clinical Decision Support System} -- CDSS), việc dựa đơn thuần vào một mô hình học sâu hộp đen (\emph{black-box Deep Learning}) hoặc một bộ câu hỏi cố định đều bộc lộ hạn chế: mô hình học sâu thiếu khả năng trích dẫn phác đồ y khoa chính thống, trong khi hệ thống tra cứu văn bản thuần túy không định lượng được xác suất nguy cơ đa biến khi dữ liệu bệnh nhân bị khuyết thiếu.

Trong giai đoạn phát triển trước của dự án, các thành phần nền tảng sau đây đã được xây dựng và kiểm nghiệm độc lập:
\begin{enumerate}
  \item \textbf{Tổ hợp mô hình Học sâu và Học máy (\emph{Deep Learning \& Machine Learning Models}) dự đoán đái tháo đường}: Được huấn luyện và đánh giá có hệ thống thông qua chuỗi 10~notebook chuẩn hóa (\texttt{00\_data\_preprocessing.ipynb} đến \texttt{09\_model\_evaluation.ipynb}) trên dữ liệu dịch tễ thực tế từ Khảo sát Sức khỏe và Dinh dưỡng Quốc gia Hoa Kỳ (CDC NHANES~\cite{nhanes}), Hệ thống Giám sát Yếu tố Nguy cơ Hành vi (CDC BRFSS 2021--2023) và bộ dữ liệu nhân trắc học người Việt, xuất ra các trọng số mô hình tại \texttt{/notebook/model/} (bao gồm mạng nơ-ron sâu có kết nối tắt \texttt{nhanes\_mlp.pt}, cây tăng cường gradient \texttt{nhanes\_gbdt.joblib}, mạng TabNet \texttt{tabnet\_attention.pt}, mạng tích chập ResNet-18 phân độ võng mạc \texttt{cnn\_retinopathy\_resnet18.pt} và mạng Bi-LSTM dự báo chuỗi glucose \texttt{rnn\_glucose\_lstm.pt}).
  \item \textbf{Hệ thống Đồ thị Tri thức (\emph{Knowledge Graph}) và Truy xuất Tăng cường Sinh văn bản (\emph{RAG}) trên Neo4j và Qdrant}: Khởi tạo cơ sở dữ liệu đồ thị Neo4j~\cite{neo4j} lưu trữ mạng lưới thực thể y khoa (bệnh lý, yếu tố nguy cơ, xét nghiệm, biến chứng, thuốc điều trị và can thiệp lối sống) kết hợp chỉ mục vector ngữ nghĩa Qdrant~\cite{qdrant} và tìm kiếm toàn văn SQLite FTS5~\cite{sqlitefts5}.
  \item \textbf{Kho ngữ liệu y khoa chuẩn hóa đa nguồn}: Kết hợp Khuyến cáo \emph{ADA Standards of Care 2024}~\cite{ada2024}, Quyết định số 5481/QĐ-BYT của Bộ Y tế Việt Nam~\cite{byt5481}, phác đồ điều trị cá thể hóa từ Hệ thống Y tế Vinmec và dữ liệu dược điển thực tế từ Hệ thống Nhà thuốc FPT Long Châu.
\end{enumerate}

Bước phát triển trọng tâm của báo cáo này là tích hợp toàn bộ các thành phần rời rạc nói trên thành một \textbf{Hệ thống Đa Tác tử Lai (\emph{Hybrid Multiagent System} -- Hybrid MAS)} thống nhất~\cite{wooldridge2009,guo2024mas}. Hệ thống đề xuất kết hợp hài hòa sáu mô thức Trí tuệ Nhân tạo (\emph{AI Paradigms}):
\begin{itemize}
  \item \textbf{Học sâu và Học máy kết hợp (\emph{Deep Learning \& Calibrated Ensemble ML})}: Thực hiện suy diễn xác suất nguy cơ mắc bệnh, phân tích ảnh đáy mắt và dự báo chuỗi đường huyết liên tục (CGM) với cơ chế chịu lỗi khi dữ liệu đầu vào khuyết thiếu (\emph{Masked-Feature Augmentation}).
  \item \textbf{Đồ thị Tri thức và RAG lai (\emph{Knowledge Graph \& Hybrid RAG})}: Truy xuất ngữ cảnh lâm sàng có cấu trúc từ Neo4j và các đoạn hướng dẫn điều trị từ Qdrant/FTS5 theo cơ chế gộp xếp hạng nghịch đảo (\emph{Reciprocal Rank Fusion} -- RRF~\cite{cormack2009}).
  \item \textbf{Suy luận dựa trên luật và Phản biện đồng thuận (\emph{Rule-Based Reasoning \& Multiagent Debate})}: Áp dụng tường minh các tiêu chuẩn chẩn đoán của ADA 2024 / Bộ Y tế, kiểm tra phạm vi sàng lọc và tổ chức tranh biện đa vòng để phát hiện mâu thuẫn giữa mô hình học sâu và xét nghiệm sinh hóa~\cite{du2024debate}.
  \item \textbf{Mô hình Ngôn ngữ Lớn cục bộ (\emph{Local Large Language Models} -- Ollama / Qwen2.5)}: Thực hiện bóc tách thực thể lâm sàng từ văn bản tự nhiên / bệnh án đa phương thức và sinh báo cáo giải thích cơ chế bệnh sinh bằng tiếng Việt có neo chặt vào bằng chứng (\emph{grounded generation}).
  \item \textbf{Học tăng cường (\emph{Reinforcement Learning})}: Học chính sách gợi ý lộ trình can thiệp bậc thang bằng thuật toán Tabular Q-Learning~\cite{watkins1992,sutton2018} trong môi trường mô phỏng có kiểm soát bởi Cổng an toàn lâm sàng (\emph{Clinical Safety Gate}).
  \item \textbf{Điều phối Đa tác tử (\emph{Multiagent Coordination})}: Quản lý luồng thực thi, phân bổ tác vụ bất đồng bộ và tổng hợp đầu ra thành báo cáo hỗ trợ quyết định lâm sàng thống nhất.
\end{itemize}

\section{Mục tiêu dự án (Project Objectives)}
\label{sec:project-objectives}

Bám sát đề cương kỹ thuật của học phần Phát triển Hệ thống Thông minh, dự án đặt ra tám mục tiêu cụ thể sau đây:
\begin{enumerate}
  \item \textbf{Thiết kế kiến trúc đa tác tử mô-đun (\emph{Modular Multiagent Architecture})}: Phân tách hệ thống thành 6~tầng chức năng rõ ràng và 8~tác tử độc lập, giao tiếp qua bus thông điệp chuẩn hóa mà không phụ thuộc cứng vào nhau.
  \item \textbf{Tái sử dụng hiệu quả các mô hình Học sâu và Học máy đã huấn luyện}: Đóng gói các mô hình trong \texttt{/notebook/model/} (\texttt{nhanes\_mlp.pt}, \texttt{nhanes\_gbdt.joblib}, \texttt{brfss\_mlp.pt}, \texttt{brfss\_gbdt.joblib}, \texttt{cnn\_retinopathy\_resnet18.pt}, \texttt{rnn\_glucose\_lstm.pt}) thành dịch vụ suy diễn chuẩn hóa trong \agent{PredictionAgent} mà không cần huấn luyện lại khi vận hành.
  \item \textbf{Tái sử dụng và mở rộng Đồ thị Tri thức Neo4j}: Tích hợp trực tiếp đồ thị Neo4j (43~nút, 55~cạnh, bao gồm 12~nút thuốc \texttt{:Drug} từ ADA, Bộ Y tế, Vinmec và Nhà thuốc Long Châu) vào luồng suy luận của \agent{KnowledgeAgent}.
  \item \textbf{Định nghĩa tường minh trách nhiệm và hợp đồng giao diện của từng tác tử}: Chuẩn hóa lược đồ đầu vào (\emph{Input Schema}), đầu ra (\emph{Output Schema}) bằng Pydantic và bảng giao diện hàm cho mọi tác tử trong \texttt{/src/agents/}.
  \item \textbf{Hiện thực hóa cơ chế giao tiếp và điều phối tác tử}: Xây dựng \texttt{MessageBus} bất đồng bộ hỗ trợ 5~hành vi ngôn trung FIPA-ACL (\texttt{REQUEST}, \texttt{INFORM}, \texttt{CHALLENGE}, \texttt{RESPONSE}, \texttt{FAILURE}) dưới sự điều phối của \agent{Orchestrator}.
  \item \textbf{Xây dựng báo cáo hỗ trợ quyết định lâm sàng có cấu trúc (\emph{Structured Decision-Support Report})}: Tổng hợp báo cáo 5~phần chuyên sâu gồm: (1)~Phân tầng nguy cơ \& đối chiếu xét nghiệm ADA/BYT, (2)~Suy luận ML/DL \& phản biện hội đồng, (3)~Chuỗi suy luận đồ thị Neo4j GraphRAG \& tầm soát biến chứng, (4)~Định hướng lối sống \& thuốc tham chiếu RL/Vinmec/Long Châu, và (5)~Biện luận cơ chế bệnh sinh cá thể hóa.
  \item \textbf{Đánh giá thực nghiệm từng tác tử và toàn bộ hệ thống MAS}: Kiểm chứng bằng bảng xếp hạng thực nghiệm 9~mô hình ML/DL và bộ kiểm thử tự động 46~kịch bản (\texttt{pytest}) bao phủ toàn bộ các nhánh luồng.
  \item \textbf{Tạo nền tảng mở rộng cho nghiên cứu hệ thống thông minh đa tác tử}: Đảm bảo khả năng thay thế mô hình LLM (Ollama \texttt{qwen2.5} / HuggingFace Qwen cục bộ), bổ sung tác tử chuyên khoa mới hoặc mở rộng cơ sở tri thức mà không làm thay đổi kiến trúc lõi.
\end{enumerate}

\chapter{Kiến trúc tổng thể hệ thống}
\label{chap:architecture}

\section{Kiến trúc phân tầng sáu lớp (Layered Architecture)}
\label{sec:layered-architecture}

Hệ thống Đa Tác tử Lai hỗ trợ quyết định lâm sàng bệnh đái tháo đường được tổ chức thành sáu tầng chức năng (\emph{six-layer architecture}), đảm bảo sự tách biệt rõ ràng giữa giao diện người dùng, xử lý dữ liệu đầu vào, điều phối tác tử, suy diễn thông minh, lưu trữ mô hình -- tri thức và kiểm toán vết thực thi:

\begin{enumerate}[label=\textbf{Tầng \arabic*.}, wide=0pt, leftmargin=1.8cm]
  \item \textbf{Tầng Tương tác Người dùng (\emph{User Interaction Layer})}:
  Cung cấp giao diện Web tương tác trực quan (xây dựng bằng React 18 và Vite tại thư mục \texttt{/web}) phục vụ ba nhóm đối tượng: người bệnh cần sàng lọc nguy cơ sớm, bác sĩ cần đối chiếu nhanh bằng chứng phác đồ, và sinh viên / nhà nghiên cứu cần theo dõi luồng suy luận đa tác tử. Tầng này hỗ trợ nhập hội thoại tự nhiên tiếng Việt, tải lên bệnh án đa phương thức (PDF, ảnh chụp phiếu xét nghiệm, DOCX, CSV), ghi âm giọng nói và hiển thị trực tiếp tiến trình suy luận thời gian thực qua giao thức \emph{Server-Sent Events} (SSE).

  \item \textbf{Tầng Cổng API và Xử lý Dữ liệu (\emph{API and Data Processing Layer})}:
  Được hiện thực hóa bằng FastAPI~\cite{fastapi} tại \texttt{src/api/} (\texttt{routes\_chat.py}, \texttt{routes\_session.py}, \texttt{routes\_system.py}) kết hợp mô-đun tiền xử lý \texttt{src/document\_processing/} và \texttt{src/extraction/}. Tầng này tiếp nhận yêu cầu HTTP/SSE, lưu trữ tệp bệnh án gốc lên hệ thống lưu trữ đối tượng SeaweedFS S3 (\texttt{s3\_bucket\_uploads}), phân loại định dạng tài liệu, chuẩn hóa đơn vị đo lường sinh hóa (ví dụ: quy đổi tự động glucose từ $\text{mmol/L}$ sang $\text{mg/dL}$ với hệ số $18{,}0182$, creatinin từ $\mu\text{mol/L}$ sang $\text{mg/dL}$, tính chỉ số BMI và mức lọc cầu thận ước tính eGFR) và kiểm định tính hợp lệ của lược đồ \texttt{PatientProfile}.

  \item \textbf{Tầng Điều phối Đa tác tử (\emph{Multiagent Coordination Layer})}:
  Bao gồm tác tử điều phối trung tâm \agent{Orchestrator} (\texttt{src/agents/orchestrator.py}) và kênh truyền thông bất đồng bộ \texttt{MessageBus} (\texttt{src/communication/message\_bus.py}). Tầng này chịu trách nhiệm định tuyến thông điệp theo chuẩn FIPA-ACL~\cite{fipa2002}, phân bổ tác vụ chạy tuần tự hoặc song song (\texttt{asyncio.gather}), điều khiển vòng lặp chất vấn -- phản biện đa vòng (\emph{multi-round debate}) và hợp nhất kết quả từ các tác tử chuyên môn.

  \item \textbf{Tầng Tác tử Thông minh (\emph{Intelligent Agent Layer})}:
  Chứa các tác tử chuyên biệt đảm nhiệm từng mảng suy luận độc lập: tác tử dữ liệu và hội thoại bệnh nhân (\agent{ConversationAgent}, \agent{DocumentAgent}), tác tử kiểm tra phạm vi và làm rõ thông tin (\agent{ClarificationAgent}), tác tử dự đoán Học sâu và Học máy (\agent{PredictionAgent}), tác tử truy xuất Đồ thị Tri thức và RAG (\agent{KnowledgeAgent}), tác tử phản biện luật và đồng thuận (\agent{ConsensusCriticAgent}), tác tử Học tăng cường (\agent{TreatmentRLAgent}) và tác tử giải thích lâm sàng (\agent{ResponseAgent}).

  \item \textbf{Tầng Tri thức và Mô hình (\emph{Knowledge and Model Layer})}:
  Bao gồm kho trọng số mô hình Học sâu / Học máy đã huấn luyện tại \texttt{/notebook/model/}, cơ sở dữ liệu đồ thị Neo4j~\cite{neo4j} lưu trữ đồ thị tri thức đái tháo đường, cơ sở dữ liệu vector Qdrant~\cite{qdrant} kết hợp chỉ mục toàn văn SQLite FTS5~\cite{sqlitefts5} (nạp từ \texttt{data/output/kb/kb\_sources/}), mô hình ngôn ngữ lớn cục bộ Ollama / Qwen2.5 (\texttt{qwen2.5:3b}, \texttt{Qwen2.5-1.5B-Instruct}) cùng môi trường mô phỏng MDP cho Học tăng cường.

  \item \textbf{Tầng Đánh giá và Kiểm toán (\emph{Evaluation and Audit Layer})}:
  Được thực thi bởi bộ thu thập vết \texttt{TraceCollector} (\texttt{src/communication/trace.py}) và cơ sở dữ liệu quan hệ SQLite/PostgreSQL (\texttt{src/storage/models.py}). Tầng này ghi nhận toàn bộ chuỗi thông điệp trao đổi giữa các tác tử, thời gian trễ từng bước (ms), phiên bản mô hình, độ tin cậy, nguồn trích dẫn và lưu trữ lịch sử hội chẩn (\texttt{Consultation}) phục vụ kiểm toán lâm sàng và đánh giá thực nghiệm.
\end{enumerate}

\section{Sơ đồ kiến trúc tổng thể (Architecture Diagram)}
\label{sec:architecture-diagram}

\figref{fig:mas-architecture} minh họa kiến trúc tổng thể của Hệ thống Đa Tác tử Lai hỗ trợ quyết định lâm sàng bệnh đái tháo đường, thể hiện đầy đủ luồng đi từ người dùng qua Cổng API, Tác tử Điều phối đến các Tác tử Chuyên môn, Tầng Mô hình -- Tri thức và Báo cáo Hỗ trợ Quyết định đầu ra.

\begin{figure}[H]
\centering
\begin{adjustbox}{max width=\linewidth}
\begin{tikzpicture}[
  font=\footnotesize,
  node distance=0.55cm and 0.45cm,
  box/.style={draw=black!75, rounded corners=3pt, align=center, minimum height=0.95cm, inner sep=4pt, fill=white, line width=0.6pt},
  layer/.style={draw=blue!65!black, fill=blue!6, rounded corners=3pt, align=center, minimum width=6.2cm, minimum height=0.92cm, inner sep=4pt, line width=0.7pt},
  agentbox/.style={draw=teal!70!black, fill=teal!6, rounded corners=3pt, align=center, minimum width=3.15cm, minimum height=1.25cm, inner sep=3.5pt, line width=0.65pt},
  storebox/.style={draw=orange!75!black, fill=orange!8, rounded corners=3pt, align=center, minimum width=3.5cm, minimum height=1.05cm, inner sep=3.5pt, line width=0.65pt},
  arr/.style={-{Stealth[length=2.2mm,width=1.6mm]}, line width=0.6pt, draw=black!75}
]

% Layer 1 & 2
\node[layer] (ui) {\textbf{Người bệnh / Bác sĩ / Sinh viên}\\Giao diện Web tương tác (\texttt{web/} -- React \& SSE)};
\node[layer, below=0.45cm of ui] (api) {\textbf{API Gateway \& Xử lý Dữ liệu (\texttt{src/api/})}\\Kiểm định đầu vào, chuẩn hóa đơn vị và bóc tách yêu cầu};
\node[layer, below=0.45cm of api, fill=blue!12, draw=blue!80!black] (coord) {\textbf{Coordinator Agent (\agent{Orchestrator})}\\Điều phối tác vụ qua \texttt{MessageBus} và tổng hợp kết quả};

% Layer 4: Specialized Agents (Row of 5 core agents)
\node[agentbox, below=0.75cm of coord] (rule) {\textbf{Rule-Based \& Critic}\\\agent{ClarificationAgent}\\\agent{ConsensusCritic}};
\node[agentbox, left=0.32cm of rule] (dl) {\textbf{DL Prediction Agent}\\\agent{PredictionAgent}\\(MLP + GBDT + CNN/RNN)};
\node[agentbox, left=0.32cm of dl] (data) {\textbf{Patient Data Agent}\\\agent{ConversationAgent}\\\agent{DocumentAgent}};
\node[agentbox, right=0.32cm of rule] (kg) {\textbf{KG / RAG Agent}\\\agent{KnowledgeAgent}\\(Neo4j + Qdrant RRF)};
\node[agentbox, right=0.32cm of kg] (exp) {\textbf{Explanation Agent}\\\agent{ResponseAgent}\\(Ollama/Qwen Grounded)};

% Monitoring & RL Agents on sides of Coordinator
\node[agentbox, left=1.2cm of coord, minimum width=3.3cm] (mon) {\textbf{Monitoring / Audit}\\\texttt{TraceCollector}\\Lưu vết \& Kiểm toán};
\node[agentbox, right=1.2cm of coord, minimum width=3.3cm] (rl) {\textbf{RL Agent}\\\agent{TreatmentRLAgent}\\Môi trường MDP + Safety Gate};

% Layer 5: Knowledge & Model Stores
\node[storebox, below=0.75cm of dl] (models) {\textbf{Mô hình DL/ML đã huấn luyện}\\\texttt{nhanes\_mlp.pt}, \texttt{nhanes\_gbdt.joblib}\\\texttt{cnn\_retinopathy}, \texttt{rnn\_glucose}};
\node[storebox, below=0.75cm of kg] (neo4j) {\textbf{Đồ thị Tri thức Neo4j \& Qdrant}\\43 nút, 55 cạnh (12 nút \texttt{:Drug})\\ADA 2024, BYT, Vinmec, Long Châu};
\node[storebox, below=0.75cm of exp] (report) {\textbf{Báo cáo Hỗ trợ Quyết định}\\Dự báo + Bằng chứng + Phác đồ\\+ Giải thích cơ chế + Giới hạn};

% Connections
\draw[arr] (ui) -- (api);
\draw[arr] (api) -- (coord);
\draw[arr] (coord) -- (mon);
\draw[arr] (coord) -- (rl);

\draw[arr] (coord.south) -- ++(0,-0.32cm) -| (data.north);
\draw[arr] (coord.south) -- ++(0,-0.32cm) -| (dl.north);
\draw[arr] (coord.south) -- (rule.north);
\draw[arr] (coord.south) -- ++(0,-0.32cm) -| (kg.north);
\draw[arr] (coord.south) -- ++(0,-0.32cm) -| (exp.north);

\draw[arr] (dl.south) -- (models.north);
\draw[arr] (kg.south) -- (neo4j.north);
\draw[arr] (exp.south) -- (report.north);

\end{tikzpicture}
\end{adjustbox}
\caption{Kiến trúc Hệ thống Đa Tác tử Lai (Hybrid MAS) hỗ trợ quyết định lâm sàng bệnh đái tháo đường.}
\label{fig:mas-architecture}
\end{figure}

\section{Giải thích kiến trúc (Architecture Interpretation)}
\label{sec:architecture-interpretation}

Trong kiến trúc trên (\figref{fig:mas-architecture}), mỗi thành phần đóng một vai trò chuyên biệt nhưng phối hợp chặt chẽ trong một chu trình khép kín:
\begin{itemize}
  \item \textbf{Tác tử Điều phối (\emph{Coordinator Agent} -- \agent{Orchestrator})} là hạt nhân điều khiển trung tâm của toàn hệ thống. Cần nhấn mạnh rằng \agent{Orchestrator} không tự mình thực hiện suy diễn học sâu hay tra cứu y văn, tức là không thay thế trí thông minh của các tác tử chuyên môn. Thay vào đó, nó quản lý vòng đời phiên hội chẩn, điều khiển thứ tự kích hoạt các tác tử chuyên biệt, phát hiện các điểm bất đồng để khởi tạo vòng chất vấn (\texttt{CHALLENGE}) và tích hợp các đầu ra thành một cấu trúc dữ liệu thống nhất.
  \item \textbf{Tác tử Dữ liệu Bệnh nhân (\emph{Patient Data Agent} -- tổ hợp \agent{ConversationAgent} và \agent{DocumentAgent})} chịu trách nhiệm tiếp nhận dữ liệu thô từ tin nhắn hoặc tệp đính kèm, trích xuất các chỉ số sinh hóa và yếu tố nguy cơ, chuẩn hóa đơn vị đo lường và kiểm tra mức độ đầy đủ của hồ sơ (\texttt{PatientProfile.is\_ready()}).
  \item \textbf{Tác tử Dự đoán Học sâu (\emph{DL Prediction Agent} -- \agent{PredictionAgent})} chịu trách nhiệm thực thi suy diễn bằng các mô hình Học sâu (\texttt{MLPRisk} PyTorch) và Học máy (\texttt{GradientBoostingRisk}) đã huấn luyện sẵn trên tập dữ liệu lâm sàng NHANES và tập hành vi quy mô lớn BRFSS, đồng thời kích hoạt mô hình thị giác ResNet-18 hoặc chuỗi thời gian Bi-LSTM khi có dữ liệu tương ứng.
  \item \textbf{Tác tử Truy xuất Tri thức (\emph{KG/RAG Agent} -- \agent{KnowledgeAgent})} truy hồi thông tin chuyên môn từ Đồ thị Tri thức Neo4j và kho văn bản Qdrant/FTS5, cung cấp các đường đi nhân quả (\texttt{graph\_paths}) và các đoạn hướng dẫn lâm sàng có gắn nguồn trích dẫn rõ ràng.
  \item \textbf{Tác tử Luật và Phản biện (\emph{Rule-Based \& Critic Agent} -- tổ hợp \agent{ClarificationAgent} và \agent{ConsensusCriticAgent})} thực thi các luật y khoa tường minh đã được kiểm chứng (tiêu chuẩn ADA 2024, Quyết định 5481/QĐ-BYT, quy tắc loại trừ thai kỳ / nhi khoa / bệnh cấp tính) và đối chiếu chéo giữa xác suất mô hình với kết quả xét nghiệm để gắn cờ cảnh báo (\emph{review flags}).
  \item \textbf{Tác tử Giải thích (\emph{Explanation Agent} -- \agent{ResponseAgent})} xây dựng báo cáo tư vấn có cấu trúc 5~phần từ đầu ra của các tác tử khác, sử dụng mô hình ngôn ngữ lớn cục bộ (Ollama / Qwen2.5) kết hợp bộ biện luận sinh lý bệnh để giải thích rõ ràng, không bịa đặt thông tin ngoài bằng chứng.
  \item \textbf{Tác tử Giám sát (\emph{Monitoring Agent} -- \texttt{TraceCollector})} ghi lại nhật ký kiểm toán (\emph{audit trail}) cho từng thông điệp, đảm bảo khả năng truy vết đầy đủ từ câu trả lời cuối cùng ngược về từng bước suy diễn.
  \item \textbf{Tác tử Học tăng cường (\emph{RL Agent} -- \agent{TreatmentRLAgent})} đóng vai trò mô-đun mở rộng nâng cao, vận hành trên bảng giá trị Q-table đã học trong môi trường mô phỏng kết hợp Cổng an toàn lâm sàng để gợi ý lộ trình can thiệp lối sống và nhóm thuốc tham chiếu (ADA / Vinmec / Nhà thuốc Long Châu) phục vụ trao đổi với bác sĩ.
\end{itemize}

\chapter{Phân loại tác tử và trách nhiệm}
\label{chap:agents}

\section{Phân loại tác tử (Agent Classification)}
\label{sec:agent-classification}

Hệ thống Đa Tác tử Lai được thiết kế theo mô hình tổ chức chuyên môn hóa: mỗi tác tử đại diện cho một vai trò chức năng trong hội đồng hỗ trợ quyết định y khoa và vận dụng mô thức Trí tuệ Nhân tạo phù hợp nhất với nhiệm vụ của mình. \tabref{tab:agent-classification} trình bày bảng phân loại đầy đủ các tác tử theo đề cương kiến trúc chuẩn, ánh xạ trực tiếp sang các lớp (\emph{class}) thực thi trong thư mục \texttt{src/agents/} và \texttt{src/communication/}.

\begin{table}[htbp]
\centering
\caption{Phân loại các tác tử trong hệ thống, mô thức AI và trách nhiệm chính.}
\label{tab:agent-classification}
\footnotesize
\renewcommand{\arraystretch}{1.1}
\begin{tabularx}{\linewidth}{@{}L{3.7cm}L{3.2cm}Y@{}}
\toprule
\textbf{Tác tử (Class thực thi)} & \textbf{Mô thức AI (AI Paradigm)} & \textbf{Trách nhiệm chính (Main Responsibility)} \\
\midrule
\textbf{Patient Data Agent}\newline (\agent{ConversationAgent}, \agent{DocumentAgent}) &
Xử lý tất định kết hợp trích xuất cấu trúc (\emph{Deterministic \& Structured Extraction}) &
Phân loại ý định, bóc tách chỉ số lâm sàng từ hội thoại hoặc bệnh án (PDF, ảnh, DOCX, CSV), chuẩn hóa đơn vị đo lường, kiểm tra kiểu dữ liệu, giá trị khuyết thiếu và tính nhất quán của \texttt{PatientProfile}. \\
\addlinespace[3pt]
\textbf{DL Prediction Agent}\newline (\agent{PredictionAgent}) &
Học sâu \& Học máy kết hợp (\emph{Deep Learning \& Calibrated Ensemble ML}) &
Thực thi mô hình Học sâu (\texttt{nhanes\_mlp.pt}) và Học máy (\texttt{nhanes\_gbdt.joblib}) đã huấn luyện sẵn, đối chiếu với quần thể BRFSS (\texttt{brfss\_mlp.pt}, \texttt{brfss\_gbdt.joblib}), phân tích chuỗi CGM (\texttt{rnn\_glucose\_lstm.pt}) và ảnh võng mạc (\texttt{cnn\_retinopathy\_resnet18.pt}). \\
\addlinespace[3pt]
\textbf{Rule-Based Agent}\newline (\agent{ClarificationAgent}, \agent{ConsensusCriticAgent}) &
Suy luận biểu tượng \& Phản biện đồng thuận (\emph{Symbolic Reasoning \& Debate}) &
Áp dụng tường minh tiêu chuẩn ADA 2024 / Bộ Y tế, kiểm tra phạm vi sàng lọc, phát hiện mâu thuẫn giữa dự báo học sâu và xét nghiệm sinh hóa, tạo các cờ cảnh báo kiểm toán (\emph{auditable review flags}). \\
\addlinespace[3pt]
\textbf{KG/RAG Agent}\newline (\agent{KnowledgeAgent}) &
Đồ thị Tri thức \& RAG lai (\emph{Knowledge Graph \& Hybrid RAG}) &
Liên kết thực thể y khoa, truy vấn lân cận $k$-bước trên đồ thị Neo4j kết hợp tìm kiếm lai vector Qdrant + toàn văn FTS5 (RRF), trả về đường đi suy luận và nguồn bằng chứng gốc. \\
\addlinespace[3pt]
\textbf{Coordinator Agent}\newline (\agent{Orchestrator}) &
Điều phối đa tác tử (\emph{Multiagent Orchestration}) &
Lập lịch kích hoạt các tác tử chuyên môn qua \texttt{MessageBus}, điều khiển vòng lặp tranh biện, hợp nhất đầu ra và xử lý lỗi dừng an toàn. \\
\addlinespace[3pt]
\textbf{Explanation Agent}\newline (\agent{ResponseAgent}) &
Suy luận bằng LLM cục bộ kết hợp khung biện luận (\emph{LLM \& Grounded Reasoning}) &
Tổng hợp báo cáo lâm sàng 5~phần có dẫn chứng bằng tiếng Việt sử dụng Ollama / Qwen2.5 và cơ chế biện luận sinh lý bệnh, tuyệt đối không sinh dữ liệu ngoài bằng chứng. \\
\addlinespace[3pt]
\textbf{Monitoring Agent}\newline (\texttt{TraceCollector}) &
Giám sát \& Kiểm toán (\emph{Monitoring \& Audit}) &
Theo dõi thời gian thực từng bước thực thi, lưu vết thông điệp, độ trễ (ms), phiên bản mô hình và bản ghi kiểm toán vào cơ sở dữ liệu. \\
\addlinespace[3pt]
\textbf{RL Agent}\newline (\agent{TreatmentRLAgent}) &
Học tăng cường có cổng an toàn (\emph{Safe Reinforcement Learning}) &
Học và tra cứu chính sách can thiệp tối ưu trên bảng Q-table (\texttt{rl\_qtable.npz}) trong không gian mô phỏng 216~trạng thái, kết hợp Cổng an toàn lâm sàng và thuốc tham chiếu Vinmec / Long Châu. \\
\bottomrule
\end{tabularx}
\end{table}

\section{Nguyên lý thiết kế tác tử (Agent Design Principles)}
\label{sec:agent-design-principles}

Để đảm bảo tính tin cậy của một hệ thống hỗ trợ quyết định trong lĩnh vực y sinh, mọi tác tử trong dự án đều tuân thủ nghiêm ngặt sáu nguyên lý thiết kế kỹ thuật phần mềm và tác tử thông minh:

\begin{enumerate}
  \item \textbf{Nguyên lý đơn trách nhiệm (\emph{Single Responsibility Principle})}:
  Mỗi tác tử chỉ đảm nhận một miền nhiệm vụ xác định rõ ràng và không chồng lấn. Cụ thể, \agent{PredictionAgent} chỉ tính toán xác suất thống kê từ trọng số mạng nơ-ron/cây quyết định mà không phán xét tiêu chuẩn chẩn đoán; ngược lại, \agent{KnowledgeAgent} chỉ đối chiếu ngưỡng sinh hóa và truy xuất y văn từ Neo4j/Qdrant mà không can thiệp vào quá trình suy diễn của mô hình học sâu.

  \item \textbf{Tính mô-đun và cô lập (\emph{Modularity})}:
  Mỗi tác tử kế thừa từ lớp cơ sở trừu tượng \texttt{BaseAgent} (\texttt{src/communication/message\_bus.py}), khai báo tập thông điệp tiếp nhận qua thuộc tính \texttt{handles} và có thể được khởi tạo, kiểm thử đơn vị (\emph{unit test}) hoàn toàn độc lập với các tác tử còn lại.

  \item \textbf{Giao diện và hợp đồng dữ liệu tường minh (\emph{Explicit Interfaces})}:
  Các tác tử không gọi trực tiếp phương thức nội bộ của nhau mà chỉ trao đổi qua đối tượng \texttt{AgentMessage} bất biến (\texttt{frozen=True}) và lược đồ dữ liệu chuẩn hóa \texttt{PatientProfile} (Pydantic). Mọi giá trị lâm sàng đều đi kèm độ tin cậy (\texttt{confidence}), nguồn trích xuất (\texttt{source}), đoạn trích nguyên văn (\texttt{evidence}) và cờ cảnh báo đơn vị (\texttt{flags}).

  \item \textbf{Khả năng truy vết toàn diện (\emph{Traceability})}:
  Mọi kết luận quan trọng trong báo cáo cuối cùng đều gắn liền với bằng chứng gốc (\emph{evidence}) hoặc lý giải định lượng (\emph{rationale}): dự báo học sâu đi kèm độ đóng góp của từng đặc trưng (\texttt{top\_drivers}); khuyến cáo lâm sàng đi kèm đường đi trên đồ thị Neo4j (\texttt{graph\_paths}) và mã định danh đoạn văn (\texttt{chunk\_ids}) trỏ về tài liệu gốc của ADA 2024, Bộ Y tế, Vinmec và Nhà thuốc FPT Long Châu.

  \item \textbf{Nhận thức lỗi và suy giảm êm (\emph{Failure Awareness \& Graceful Degradation})}:
  Các lỗi định dạng tệp, thiếu chỉ số bắt buộc hoặc tình huống nằm ngoài phạm vi sàng lọc (mang thai, nhi khoa, cấp cứu tăng đường huyết) đều được trả về tường minh thông qua thông điệp \texttt{FAILURE} hoặc câu hỏi làm rõ chặn luồng (\texttt{blocking=True}) thay vì bị che giấu. Khi dịch vụ LLM cục bộ tạm ngừng, từng tác tử tự động chuyển sang nhánh xử lý tất định (\emph{deterministic fallback}) mà không làm gián đoạn hệ thống.

  \item \textbf{Khả năng mở rộng (\emph{Extensibility})}:
  Nhờ cơ chế đăng ký động \texttt{bus.register(agent)} trên \texttt{MessageBus}, việc thay thế một mô hình học sâu mới trong \texttt{/notebook/model/} hoặc bổ sung một tác tử chuyên khoa mới (ví dụ: tác tử dinh dưỡng cá thể hóa) hoàn toàn không đòi hỏi thiết kế lại kiến trúc tổng thể của hệ thống.
\end{enumerate}

\chapter{Mô hình hóa toán học Hệ Đa Tác tử}
\label{chap:math}

\section{Định nghĩa hình thức Hệ Đa Tác tử (Multiagent System Definition)}
\label{sec:mas-definition}

Theo lý thuyết hệ thống đa tác tử~\cite{wooldridge2009}, toàn bộ Hệ thống Đa Tác tử Lai hỗ trợ quyết định lâm sàng bệnh đái tháo đường được mô hình hóa hình thức bởi một bộ năm (\emph{5-tuple}):
\begin{equation}
  M = (\mathcal{A}, \mathcal{S}, \mathcal{O}, \mathcal{C}, \mathcal{F}),
  \label{eq:mas-tuple}
\end{equation}
trong đó:
\begin{itemize}
  \item $\mathcal{A} = \{A_1, A_2, \dots, A_n\}$ là tập hữu hạn các tác tử chuyên biệt trong hệ thống.
  \item $\mathcal{S}$ là không gian trạng thái toàn cục của hệ thống (\emph{global system state space}), tại mỗi thời điểm hội thoại $t$ được biểu diễn bởi cặp trạng thái $s_t = (P_t, H_t) \in \mathcal{S}$, với $P_t \in \mathcal{X}_{\text{patient}}$ là hồ sơ lâm sàng tích lũy (\texttt{PatientProfile}) gồm các biến sinh hóa/nhân trắc có gắn độ tin cậy và $H_t$ là lịch sử các lượt hội chẩn trong phiên.
  \item $\mathcal{O} = \{O_1, O_2, \dots, O_n\}$ là tập các phép chiếu quan sát của từng tác tử (\emph{agent observations}), trong đó mỗi tác tử $A_i$ chỉ quan sát một tập con thông tin liên quan đến miền chuyên môn của mình thông qua hàm quan sát $o_i: \mathcal{S} \to X_i$.
  \item $\mathcal{C}$ là cơ chế truyền thông (\emph{communication mechanism}), được định nghĩa bởi không gian thông điệp có cấu trúc $\mathcal{M}_{\text{msg}}$ truyền qua kênh bất đồng bộ \texttt{MessageBus} với tập hành vi ngôn trung $\mathcal{P}_{\text{ACL}} = \{\texttt{REQUEST}, \allowbreak \texttt{INFORM}, \allowbreak \texttt{CHALLENGE}, \allowbreak \texttt{RESPONSE}, \allowbreak \texttt{FAILURE}\}$.
  \item $\mathcal{F}$ là hàm điều phối và hợp nhất (\emph{coordination and aggregation function}) do tác tử điều phối thực thi nhằm kết hợp đầu ra của các tác tử thành viên thành quyết định tổng thể của hệ thống.
\end{itemize}

Đối với hệ thống đề xuất trong dự án này, tập tác tử $\mathcal{A}$ được định nghĩa tường minh gồm tám thành phần chức năng:
\begin{equation}
  \mathcal{A} = \{A_{\text{data}}, A_{\text{DL}}, A_{\text{rule}}, A_{\text{KG}}, A_{\text{coord}}, A_{\text{exp}}, A_{\text{monitor}}, A_{\text{RL}}\},
  \label{eq:agent-set}
\end{equation}
trong đó:
\begin{itemize}
  \item $A_{\text{data}} = \{\agent{ConversationAgent}, \agent{DocumentAgent}\}$: tác tử bóc tách, chuẩn hóa và kiểm định dữ liệu bệnh nhân;
  \item $A_{\text{DL}} = \agent{PredictionAgent}$: tác tử suy diễn Học sâu và Học máy kết hợp;
  \item $A_{\text{rule}} = \{\agent{ClarificationAgent}, \agent{ConsensusCriticAgent}\}$: tác tử kiểm định luật lâm sàng và phản biện đồng thuận;
  \item $A_{\text{KG}} = \agent{KnowledgeAgent}$: tác tử truy xuất Đồ thị Tri thức Neo4j và Hybrid RAG;
  \item $A_{\text{coord}} = \agent{Orchestrator}$: tác tử điều phối trung tâm;
  \item $A_{\text{exp}} = \agent{ResponseAgent}$: tác tử giải thích và tổng hợp báo cáo lâm sàng;
  \item $A_{\text{monitor}} = \texttt{TraceCollector}$: tác tử giám sát và kiểm toán vết thực thi;
  \item $A_{\text{RL}} = \agent{TreatmentRLAgent}$: tác tử Học tăng cường tối ưu lộ trình can thiệp trong không gian mô phỏng.
\end{itemize}

\section{Ánh xạ đầu vào, đầu ra và hàm tổng hợp (Agent Input and Output)}
\label{sec:agent-io}

Mỗi tác tử $A_i \in \mathcal{A}$ được mô hình hóa toán học như một ánh xạ hàm từ không gian đầu vào $X_i$ sang không gian đầu ra có cấu trúc $Y_i$:
\begin{equation}
  A_i : X_i \longrightarrow Y_i.
  \label{eq:agent-map}
\end{equation}

Cụ thể, các ánh xạ thành phần chính trong hệ thống được xác định như sau:

\paragraph{1. Tác tử Dự đoán Học sâu và Học máy ($A_{\text{DL}}$):}
Nhận vector đặc trưng bệnh nhân $x \in X_{\text{patient}} \subset (\mathbb{R} \cup \{\text{NaN}\})^d$ và trả về cấu trúc dự báo $y_{\text{DL}} \in Y_{\text{prediction}}$:
\begin{equation}
  A_{\text{DL}} : X_{\text{patient}} \longrightarrow Y_{\text{prediction}}, \qquad
  p_{\text{ens}}(x) = \frac{1}{2}\Big(\sigma\big(a \cdot f_{\text{GBDT}}(x) + b\big) + f_{\text{MLP}}(x)\Big),
  \label{eq:dl-agent}
\end{equation}
trong đó $\sigma(a \cdot f_{\text{GBDT}}(x) + b)$ là xác suất của mô hình cây tăng cường gradient sau hiệu chỉnh Platt~\cite{platt1999}, $f_{\text{MLP}}(x)$ là xác suất đầu ra của mạng nơ-ron sâu có kết nối tắt (\emph{Residual MLP}), đi kèm tỷ số nguy cơ tương đối $\text{RR}(x) = p_{\text{ens}}(x) / \bar{p}_{\text{pop}}$ và độ lệch bất đồng giữa hai mô hình $\Delta_{\text{model}}(x) = |p_{\text{GBDT}}(x) - p_{\text{MLP}}(x)|$.

\paragraph{2. Tác tử Truy xuất Đồ thị Tri thức và RAG ($A_{\text{KG}}$):}
Biểu diễn ánh xạ từ không gian truy vấn và thực thể lâm sàng $Q$ sang không gian bằng chứng y khoa có nguồn gốc $E$:
\begin{equation}
  A_{\text{KG}} : Q \longrightarrow E, \qquad
  \text{RRF}(d) = \frac{1}{k_0 + r_{\text{vec}}(d)} + \frac{1}{k_0 + r_{\text{fts}}(d)} \quad (k_0 = 60),
  \label{eq:kg-agent}
\end{equation}
kết hợp tập đường đi lân cận $k$-bước ($k \le 2$) trên đồ thị Neo4j $G_{\text{Neo4j}} = (V, R_{\text{edge}})$ xuất phát từ tập nút thực thể $V_0 \subset V$ khớp với hồ sơ bệnh nhân.

\paragraph{3. Tác tử Luật và Phản biện Đồng thuận ($A_{\text{rule}}$):}
Đánh giá các tiêu chuẩn chẩn đoán ADA 2024 / Bộ Y tế và tính điểm nguy cơ hội đồng sau khi thực hiện các vòng chất vấn (\texttt{CHALLENGE}):
\begin{equation}
  S_{\text{consensus}} = \max\Big(s_{\text{lab}}(x),\, w_{\text{lab}} \cdot s_{\text{lab}}(x) + (1 - w_{\text{lab}}) \cdot p_{\text{ens}}(x)\Big),
  \label{eq:consensus-score}
\end{equation}
trong đó $s_{\text{lab}}(x) \in \{0{,}15;\, 0{,}55;\, 0{,}90\}$ tương ứng với ngưỡng bình thường, tiền đái tháo đường và ngưỡng đái tháo đường theo ADA~2024~\cite{ada2024}. Công thức $\max(\cdot)$ đảm bảo nguyên tắc an toàn lâm sàng: một xác suất mô hình thấp không bao giờ được phép hạ thấp cảnh báo khi xét nghiệm sinh hóa đã vượt ngưỡng bệnh lý.

\paragraph{4. Hàm điều phối và tổng hợp toàn hệ thống ($\mathcal{F}$):}
Tác tử điều phối $A_{\text{coord}}$ tích hợp đầu ra của các tác tử thành viên thông qua hàm tổng hợp:
\begin{equation}
  \mathcal{F}\big(y_{\text{DL}},\, y_{\text{rule}},\, y_{\text{KG}},\, y_{\text{RL}},\, y_{\text{monitor}}\big) = y_{\text{system}}.
  \label{eq:aggregation-func}
\end{equation}
Điểm cốt lõi trong thiết kế hàm $\mathcal{F}$ của hệ thống là \textbf{bảo toàn sự phân biệt rạch ròi} giữa bốn nguồn thông tin trong cấu trúc $y_{\text{system}}$: (i)~dự báo xác suất từ mô hình Học sâu / Học máy ($y_{\text{DL}}$), (ii)~kết quả đối chiếu luật lâm sàng và cờ phản biện ($y_{\text{rule}}$), (iii)~bằng chứng truy xuất từ Đồ thị Tri thức Neo4j và Qdrant ($y_{\text{KG}}$), và (iv)~gợi ý mô phỏng từ Học tăng cường ($y_{\text{RL}}$). Sự tách bạch này giúp bác sĩ và người dùng kiểm chứng độc lập từng nguồn suy luận.

\chapter{Luồng tương tác giữa các tác tử}
\label{chap:workflow}

\section{Luồng xử lý tuần tự và song song (Sequential \& Parallel Workflow)}
\label{sec:sequential-workflow}

Luồng tương tác giữa các tác tử trong một lượt hội chẩn (\emph{consultation turn}) tuân thủ quy trình 11~bước chặt chẽ, kết hợp kiểm định tuần tự ở đầu vào với suy diễn song song ở tầng chuyên môn:

\begin{enumerate}
  \item \textbf{Bước 1}: Người dùng gửi các chỉ số sức khỏe, tệp bệnh án đính kèm (nếu có) và câu hỏi tự nhiên qua giao diện Web.
  \item \textbf{Bước 2}: Tầng Cổng API (\texttt{routes\_chat.py}) tiếp nhận yêu cầu, tạo mã phiên và chuyển luồng điều khiển cho Tác tử Điều phối (\agent{Orchestrator}).
  \item \textbf{Bước 3}: Tác tử Dữ liệu Bệnh nhân (\agent{DocumentAgent} và \agent{ConversationAgent}) bóc tách chỉ số, chuẩn hóa đơn vị đo lường, kiểm tra tính hợp lệ của \texttt{PatientProfile} và phối hợp với \agent{ClarificationAgent} kiểm tra phạm vi sàng lọc.
  \item \textbf{Bước 4}: Nếu dữ liệu không hợp lệ, thiếu chỉ số bắt buộc (chưa có cả đường huyết/HbA1c lẫn cặp tuổi + BMI) hoặc gặp điều kiện chặn luồng (\texttt{blocking=True} như đang mang thai, bệnh nhi), Tác tử Điều phối dừng sớm luồng suy diễn và trả về thông báo yêu cầu làm rõ / bổ sung dữ liệu.
  \item \textbf{Bước 5}: Tác tử Dự đoán Học sâu (\agent{PredictionAgent}) thực thi song song cặp mô hình đã huấn luyện (\texttt{nhanes\_mlp.pt} + \texttt{nhanes\_gbdt.joblib}) cùng mô hình đối chứng quy mô lớn BRFSS (\texttt{brfss\_mlp.pt} + \texttt{brfss\_gbdt.joblib}).
  \item \textbf{Bước 6}: Tác tử Truy xuất Tri thức (\agent{KnowledgeAgent}) đối chiếu xét nghiệm với ngưỡng ADA 2024 / Bộ Y tế, đồng thời truy xuất đường đi trên Đồ thị Tri thức Neo4j và các đoạn hướng dẫn từ Qdrant/FTS5.
  \item \textbf{Bước 7}: Tác tử Luật và Phản biện (\agent{ConsensusCriticAgent}) đánh giá các luật lâm sàng tường minh; nếu phát hiện mâu thuẫn giữa dự báo Học sâu và xét nghiệm sinh hóa, tác tử này phát thông điệp chất vấn (\texttt{CHALLENGE}) để các tác tử liên quan giải trình.
  \item \textbf{Bước 8}: Tác tử Học tăng cường (\agent{TreatmentRLAgent}) ánh xạ trạng thái bệnh nhân vào bảng Q-table (\texttt{rl\_qtable.npz}) và áp dụng Cổng an toàn lâm sàng để chọn khuyến cáo can thiệp kèm danh mục thuốc tham chiếu.
  \item \textbf{Bước 9}: Tác tử Điều phối (\agent{Orchestrator}) thu thập và hợp nhất toàn bộ đầu ra từ các tác tử chuyên môn.
  \item \textbf{Bước 10}: Tác tử Giải thích (\agent{ResponseAgent}) sinh báo cáo hỗ trợ quyết định lâm sàng có cấu trúc 5~phần, trong khi Tác tử Giám sát (\texttt{TraceCollector}) ghi nhận toàn bộ siêu dữ liệu thực thi (\emph{execution metadata}).
  \item \textbf{Bước 11}: Báo cáo hoàn chỉnh kèm vết suy luận và trích dẫn nguồn được trả về cho người dùng.
\end{enumerate}

\section{Biểu đồ tuần tự (Sequence Diagram)}
\label{sec:sequence-diagram}

\figref{fig:sequence-diagram} trình bày biểu đồ tuần tự (\emph{sequence diagram}) mô tả luồng trao đổi thông điệp giữa người dùng, Tác tử Điều phối và các tác tử chuyên môn trong một chu trình hội chẩn chuẩn.

\begin{figure}[H]
\centering
\begin{adjustbox}{max width=\linewidth}
\begin{tikzpicture}[
  font=\scriptsize,
  actor/.style={draw=blue!75!black, fill=blue!8, rounded corners=2pt, minimum width=1.85cm, minimum height=0.8cm, align=center, font=\scriptsize\bfseries},
  msg/.style={-{Stealth[length=2mm,width=1.4mm]}, line width=0.55pt, draw=black!80},
  ret/.style={-{Stealth[length=2mm,width=1.4mm]}, line width=0.55pt, dashed, draw=teal!75!black},
  life/.style={dashed, draw=black!40, line width=0.45pt}
]

% Actors across the top
\node[actor] (u) at (0,0) {Người dùng\\(User)};
\node[actor] (c) at (2.4,0) {Coordinator\\(\agent{Orchestrator})};
\node[actor] (d) at (4.9,0) {Data Agent\\(\agent{Conversation})};
\node[actor] (dl) at (7.4,0) {DL Agent\\(\agent{Prediction})};
\node[actor] (r) at (9.9,0) {Rule/Critic\\(\agent{Consensus})};
\node[actor] (kg) at (12.4,0) {KG/RAG Agent\\(\agent{Knowledge})};
\node[actor] (e) at (14.9,0) {Explanation\\(\agent{Response})};

% Lifelines
\foreach \n in {u,c,d,dl,r,kg,e} {
  \draw[life] (\n.south) -- ++(0,-8.6cm);
}

% Step 1: User -> Coordinator
\draw[msg] (0,-0.9) -- node[above, font=\tiny]{1. Yêu cầu khám / Câu hỏi} (2.4,-0.9);

% Step 2: Coordinator -> Data Agent
\draw[msg] (2.4,-1.55) -- node[above, font=\tiny]{2. \texttt{PROCESS\_TURN} (Kiểm định)} (4.9,-1.55);
\draw[ret] (4.9,-2.15) -- node[above, font=\tiny]{3. Hồ sơ hợp lệ (\texttt{PatientProfile})} (2.4,-2.15);

% Step 3: Coordinator -> DL Agent
\draw[msg] (2.4,-2.85) -- node[above, font=\tiny]{4. \texttt{PREDICT\_RISK}} (7.4,-2.85);
\draw[ret] (7.4,-3.45) -- node[above, font=\tiny]{5. Dự báo DL/ML ($p_{\text{ens}}, \text{RR}$)} (2.4,-3.45);

% Step 4: Coordinator -> KG/RAG Agent
\draw[msg] (2.4,-4.15) -- node[above, font=\tiny]{6. \texttt{RETRIEVE\_EVIDENCE}} (12.4,-4.15);
\draw[ret] (12.4,-4.75) -- node[above, font=\tiny]{7. Bằng chứng Neo4j KG + Qdrant RAG} (2.4,-4.75);

% Step 5: Coordinator -> Rule/Critic Agent (+ Debate + RL)
\draw[msg] (2.4,-5.45) -- node[above, font=\tiny]{8. \texttt{EVALUATE\_CONSENSUS} (+ RL)} (9.9,-5.45);
\draw[ret] (9.9,-6.10) -- node[above, font=\tiny]{9. Kết quả luật, cờ phản biện \& phác đồ} (2.4,-6.10);

% Step 6: Coordinator -> Explanation Agent
\draw[msg] (2.4,-6.85) -- node[above, font=\tiny]{10. \texttt{COMPOSE\_ASSESSMENT}} (14.9,-6.85);
\draw[ret] (14.9,-7.45) -- node[above, font=\tiny]{11. Báo cáo 5 phần có dẫn chứng} (2.4,-7.45);

% Step 7: Coordinator -> User
\draw[ret] (2.4,-8.15) -- node[above, font=\tiny]{12. Trả báo cáo + Vết kiểm toán (\texttt{Trace})} (0,-8.15);

\end{tikzpicture}
\end{adjustbox}
\caption{Biểu đồ tuần tự tương tác giữa các tác tử trong Hệ thống Đa Tác tử Lai.}
\label{fig:sequence-diagram}
\end{figure}

\section{Thuật toán điều phối của Coordinator (Coordinator Workflow)}
\label{sec:coordinator-workflow}

Tác tử điều phối \agent{Orchestrator} hiện thực hóa thuật toán điều khiển 12~bước chuẩn hóa sau đây trong phương thức \texttt{run\_turn()}:
\begin{enumerate}
  \item Tiếp nhận yêu cầu từ người dùng (kèm văn bản, danh sách tệp đính kèm và trạng thái hồ sơ hiện tại).
  \item Khởi tạo mã định danh lượt xử lý (\texttt{trace\_id}) và đối tượng giám sát \texttt{TraceCollector}.
  \item Gọi Tác tử Dữ liệu Bệnh nhân (\agent{DocumentAgent} nếu có tệp, sau đó là \agent{ConversationAgent}) để bóc tách và hợp nhất \texttt{PatientProfile}.
  \item Kiểm tra trạng thái hợp lệ, ý định hội thoại và điều kiện phạm vi qua \agent{ClarificationAgent}.
  \item Nếu hồ sơ chưa đủ chỉ số tối thiểu hoặc gặp điều kiện chặn (\texttt{blocking\_questions}), trả về ngay thông báo hướng dẫn bổ sung dữ liệu.
  \item Kích hoạt Tác tử Dự đoán Học sâu (\agent{PredictionAgent}) với thông điệp \texttt{PREDICT\_RISK}.
  \item Kích hoạt Tác tử Truy xuất Tri thức (\agent{KnowledgeAgent}) với thông điệp \texttt{RETRIEVE\_EVIDENCE}.
  \item Kích hoạt Tác tử Luật và Phản biện (\agent{ConsensusCriticAgent}) với thông điệp \texttt{EVALUATE\_CONSENSUS} (bao gồm việc thực thi các thông điệp \texttt{CHALLENGE} nếu phát hiện mâu thuẫn).
  \item Kích hoạt Tác tử Học tăng cường (\agent{TreatmentRLAgent}) với thông điệp \texttt{RECOMMEND\_TREATMENT} và thu thập toàn bộ đầu ra.
  \item Gọi Tác tử Giải thích (\agent{ResponseAgent}) với thông điệp \texttt{COMPOSE\_ASSESSMENT} để tổng hợp báo cáo 5~phần.
  \item Ghi nhận toàn bộ vết thực thi (\texttt{trace.transcript()}) và biểu đồ tương tác (\texttt{trace.interaction\_graph()}).
  \item Trả về cấu trúc kết quả hoàn chỉnh cho tầng API.
\end{enumerate}

\coderef{code:orchestrator-core} trích dẫn đoạn mã thực tế của phương thức \texttt{Orchestrator.run\_turn()} trong tệp \texttt{src/agents/orchestrator.py}.

\begin{code}{Thuật toán điều phối trung tâm trong \texttt{src/agents/orchestrator.py}}
\label{code:orchestrator-core}
\begin{minted}{python}
# Bước 5 & 6: Kích hoạt song song DL Prediction Agent và KG/RAG Agent
pred_msg, know_msg = await asyncio.gather(
    self.request("PredictionAgent", "PREDICT_RISK",
                 {"profile": profile.model_dump()}, trace),
    self.request("KnowledgeAgent", "RETRIEVE_EVIDENCE",
                 {"profile": profile.model_dump()}, trace),
)
prediction, knowledge = pred_msg.payload, know_msg.payload

# Bước 7: Rule-Based & Consensus Critic Agent kiểm định và tranh biện đa vòng
consensus = await self._run_consensus_debate(profile, prediction, knowledge, trace)

# Bước 8: RL Agent tra cứu chính sách can thiệp + Safety Gate
treat_msg = await self.request("TreatmentRLAgent", "RECOMMEND_TREATMENT", {
    "profile": profile.model_dump(),
    "consensus": consensus,
    "knowledge": knowledge,
}, trace)

# Bước 10: Explanation Agent tổng hợp báo cáo 5 phần có dẫn chứng
resp = await self.request("ResponseAgent", "COMPOSE_ASSESSMENT", {
    "profile": profile.model_dump(), "prediction": prediction,
    "knowledge": knowledge, "consensus": consensus,
    "treatment": treat_msg.payload, "questions": questions,
}, trace)
\end{minted}
\end{code}

\chapter{Giao thức truyền thông và Hợp đồng dữ liệu}
\label{chap:contracts}

\section{Cấu trúc thông điệp tác tử (Agent Message Structure)}
\label{sec:message-structure}

Mọi tương tác giữa các tác tử trong hệ thống đều được thực hiện thông qua các thông điệp có cấu trúc tường minh (\emph{structured messages}), tuân theo mô hình ngôn ngữ giao tiếp tác tử FIPA-ACL~\cite{fipa2002}. Về mặt hình thức, một thông điệp tổng quát $m_i$ được biểu diễn bởi bộ sáu:
\begin{equation}
  m_i = (\text{id},\, \text{agent},\, \text{status},\, \text{payload},\, \text{evidence},\, \text{version}),
  \label{eq:message-tuple}
\end{equation}
trong đó ý nghĩa của từng trường được quy định như sau:
\begin{itemize}
  \item \texttt{id}: Mã định danh duy nhất của thông điệp (\texttt{msg\_id}) và mã phiên truy vết (\texttt{trace\_id}).
  \item \texttt{agent}: Định danh tác tử gửi (\texttt{sender}) và tác tử nhận (\texttt{receiver}).
  \item \texttt{status} / \texttt{performative}: Trạng thái thực thi hoặc hành vi ngôn trung của thông điệp, nhận một trong năm giá trị chuẩn: \texttt{REQUEST} (yêu cầu thực thi tác vụ), \texttt{INFORM} (trả kết quả thành công), \texttt{CHALLENGE} (chất vấn khi phát hiện mâu thuẫn), \texttt{RESPONSE} (giải trình phản biện), và \texttt{FAILURE} (báo lỗi tường minh).
  \item \texttt{payload}: Nội dung dữ liệu đầu ra có cấu trúc (\texttt{Dict[str, Any]}).
  \item \texttt{evidence}: Danh sách nguồn bằng chứng y khoa trích xuất từ Neo4j/Qdrant hoặc các yếu tố giải thích (\texttt{top\_drivers}, \texttt{citations}, \texttt{graph\_paths}).
  \item \texttt{version}: Phiên bản mô hình hoặc cơ sở dữ liệu tri thức đã sử dụng (ví dụ: \texttt{nhanes\_mlp\_v2}, \texttt{nhanes\_gbdt\_v2}, \texttt{neo4j\_kg\_ada2024\_longchau}).
\end{itemize}

Trong mã nguồn \texttt{src/communication/message\_bus.py}, cấu trúc này được cài đặt dưới dạng lớp dữ liệu bất biến \texttt{AgentMessage(frozen=True)} đi kèm dấu thời gian \texttt{timestamp} và độ trễ xử lý \texttt{duration\_ms}.

\section{Ví dụ thông điệp chuẩn (Example Message)}
\label{sec:example-message}

\coderef{code:example-message} minh họa một thông điệp JSON thực tế do tác tử dự đoán \agent{PredictionAgent} trả về trên \texttt{MessageBus} sau khi thực thi cặp mô hình Học sâu (\texttt{nhanes\_mlp.pt}) và Học máy (\texttt{nhanes\_gbdt.joblib}) trên hồ sơ bệnh nhân.

\begin{code}{Ví dụ thông điệp dự báo nguy cơ từ \agent{PredictionAgent}}
\label{code:example-message}
\begin{minted}{json}
{
  "request_id": "req_98f4a21c",
  "agent": "PredictionAgent",
  "performative": "INFORM",
  "msg_type": "RISK_PREDICTION",
  "status": "ok",
  "payload": {
    "prediction": 1,
    "risk_probability": 0.382,
    "relative_risk": 1.93,
    "model_disagreement": 0.168,
    "models": {
      "ml": {"name": "nhanes_gbdt", "risk_probability": 0.298, "roc_auc": 0.8167},
      "dl": {"name": "nhanes_mlp", "risk_probability": 0.466, "roc_auc": 0.8250}
    },
    "large_cohort_opinion": {
      "dataset": "BRFSS 2021-2023 (n=120,000)",
      "relative_risk": 1.42
    }
  },
  "evidence": [
    {"feature": "Tiền sử gia đình mắc ĐTĐ", "contribution": 0.114},
    {"feature": "BMI (26.4 kg/m² >= 23 châu Á)", "contribution": 0.089},
    {"feature": "Tuổi (52 tuổi >= 35)", "contribution": 0.062}
  ],
  "model_version": "nhanes_ensemble_dl_ml_v2.0"
}
\end{minted}
\end{code}

Các con số xác suất trong \coderef{code:example-message} được hiệu chỉnh xác suất thực nghiệm (Platt Scaling và Temperature Calibration), phản ánh nguy cơ thống kê dựa trên yếu tố nhân trắc -- tiền sử và luôn được chuyển tiếp cho \agent{ConsensusCriticAgent} để đối chiếu chéo với xét nghiệm sinh hóa.

\section{Thiết kế giao diện tác tử (Agent Interface Design)}
\label{sec:agent-interfaces}

Mỗi tác tử trong hệ thống công bố phương thức giao tiếp bất đồng bộ \texttt{async def handle(msg: AgentMessage) -> AgentMessage} với hợp đồng đầu vào và đầu ra chặt chẽ (\tabref{tab:agent-interfaces}).

\begin{table}[H]
\centering
\caption{Hợp đồng giao diện đầu vào và đầu ra của các tác tử trong hệ thống.}
\label{tab:agent-interfaces}
\small
\begin{tabularx}{\linewidth}{@{}L{3.3cm}L{4.2cm}Y@{}}
\toprule
\textbf{Tác tử (Agent)} & \textbf{Đầu vào (Input Contract)} & \textbf{Đầu ra (Output Contract)} \\
\midrule
\textbf{Data Agent}\newline (\agent{ConversationAgent}, \agent{DocumentAgent}) &
Văn bản tin nhắn thô, danh sách tệp bệnh án (\texttt{document\_ids}) và \texttt{PatientProfile} hiện tại. &
\texttt{PatientProfile} đã chuẩn hóa đơn vị, ý định hội thoại (\texttt{intent}) hoặc danh sách lỗi đọc tệp / trường còn thiếu. \\
\addlinespace[3pt]
\textbf{DL Agent}\newline (\agent{PredictionAgent}) &
\texttt{PatientProfile} đã kiểm định (kèm chuỗi CGM \texttt{glucose\_history} hoặc ảnh đáy mắt nếu có). &
Xác suất nguy cơ tổng hợp, xác suất từng mô hình (\texttt{ml}, \texttt{dl}, \texttt{brfss}), đóng góp đặc trưng, phân độ võng mạc ResNet-18 và dự báo LSTM. \\
\addlinespace[3pt]
\textbf{Rule Agent}\newline (\agent{ClarificationAgent}, \agent{ConsensusCriticAgent}) &
\texttt{PatientProfile}, kết quả dự báo của DL Agent và bằng chứng từ KG/RAG Agent. &
Kết quả đối chiếu ngưỡng ADA 2024 / Bộ Y tế, câu hỏi kiểm tra phạm vi, điểm nguy cơ đồng thuận và danh sách cờ cảnh báo (\texttt{safety\_guardrails}). \\
\addlinespace[3pt]
\textbf{KG/RAG Agent}\newline (\agent{KnowledgeAgent}) &
Câu hỏi ngôn ngữ tự nhiên hoặc tập yếu tố nguy cơ / xét nghiệm của bệnh nhân. &
Đường đi trên đồ thị Neo4j (\texttt{graph\_paths}), khuyến cáo tầm soát, chống chỉ định thuốc, các đoạn văn bản xếp hạng theo RRF và nguồn trích dẫn (\texttt{citations}). \\
\addlinespace[3pt]
\textbf{Coordinator}\newline (\agent{Orchestrator}) &
Yêu cầu hội chẩn từ người dùng qua API (\texttt{run\_turn}). &
Kết quả hội chẩn tổng hợp từ toàn bộ luồng tác tử kèm cấu trúc dữ liệu chi tiết. \\
\addlinespace[3pt]
\textbf{Explanation Agent}\newline (\agent{ResponseAgent}) &
Kết quả dự báo, kết luận luật/phản biện, bằng chứng KG/RAG và khuyến cáo RL. &
Báo cáo tư vấn 5~phần có cấu trúc Markdown bằng tiếng Việt, neo chặt vào nguồn trích dẫn. \\
\addlinespace[3pt]
\textbf{Monitoring Agent}\newline (\texttt{TraceCollector}) &
Siêu dữ liệu thực thi từng bước của các tác tử trên \texttt{MessageBus}. &
Bản ghi kiểm toán (\texttt{transcript}) và đồ thị tương tác (\texttt{interaction\_graph}) có gắn thời gian trễ (ms). \\
\addlinespace[3pt]
\textbf{RL Agent}\newline (\agent{TreatmentRLAgent}) &
Trạng thái lâm sàng rời rạc hóa $s \in \{0,\dots,215\}$ và cờ an toàn thận/đường huyết. &
Hành động can thiệp tối ưu $a^*$, vector giá trị $Q(s,a)$, danh mục thuốc tham chiếu (\texttt{reference\_drugs\_vn}) và kế hoạch theo dõi. \\
\bottomrule
\end{tabularx}
\end{table}

\chapter{Tích hợp Mô hình Học sâu và Đồ thị Tri thức / RAG}
\label{chap:integration}

\section{Tích hợp Mô hình Học sâu và Học máy (Deep Learning \& ML Integration)}
\label{sec:dl-integration}

Trong kiến trúc Hệ thống Đa Tác tử Lai, các mô hình Học sâu (\emph{Deep Learning}) và Học máy (\emph{Machine Learning}) đã huấn luyện tại thư mục \texttt{/notebook/model/} được đóng gói như các dịch vụ suy diễn tái sử dụng được (\emph{reusable inference services}) bên trong \agent{PredictionAgent}:
\begin{equation}
  \text{Đặc trưng bệnh nhân (\texttt{PatientProfile})} \longrightarrow \text{Mô hình DL/ML} \longrightarrow \text{Đầu ra dự báo \& Giải thích}.
  \label{eq:dl-pipeline}
\end{equation}

Theo đúng quy chuẩn tài liệu kỹ thuật, bảy thông số tích hợp cốt lõi của tổ hợp mô hình được đặc tả chi tiết như sau:

\begin{enumerate}
  \item \textbf{Kiến trúc mô hình (\emph{Model Architecture})}:
  Hệ thống huấn luyện và đánh giá 9~kiến trúc mô hình thông qua chuỗi 10~notebook tại \texttt{/notebook/} (\texttt{00\_data\_preprocessing.ipynb} đến \texttt{09\_model\_evaluation.ipynb}):
  \begin{itemize}
    \item \emph{Mô hình Học sâu dữ liệu bảng chính (\texttt{MLPRisk} -- \texttt{05\_deep\_residual\_mlp.ipynb})}: Mạng nơ-ron sâu PyTorch~\cite{paszke2019} sử dụng nối tắt phần dư (\emph{Residual Skip-Connection}), chuẩn hóa theo lớp (\emph{LayerNorm}), hàm kích hoạt SiLU và Dropout ($p=0{,}20$). Với mỗi đặc trưng $x_j$, mạng ghép nối giá trị đã chuẩn hóa Z-score với bit mặt nạ quan sát $m_j \in \{0,1\}$ (tạo vector đầu vào $2d$ chiều) giúp mạng học trực tiếp phân phối dữ liệu khuyết thiếu.
    \item \emph{Mô hình Học sâu chú ý đặc trưng (\texttt{TabNetAttention} -- \texttt{06\_tabnet\_attention.ipynb})}: Mạng chú ý tuần tự lựa chọn đặc trưng thưa (\emph{sparse feature selection}) cho dữ liệu bảng.
    \item \emph{Mô hình Học máy đồng hành (\texttt{GradientBoostingRisk} -- \texttt{04\_hist\_gradient\_boosting.ipynb})}: Cây tăng cường gradient theo histogram (\emph{HistGradientBoosting}~\cite{friedman2001,ke2017}) có xử lý nhánh \texttt{NaN} tự nhiên (\emph{native missing-value routing}) kết hợp hiệu chỉnh xác suất Platt Scaling~\cite{platt1999} cùng các mô hình đối chứng Logistic Regression (\texttt{01}), Random Forest (\texttt{02}), Extra Trees (\texttt{03}).
    \item \emph{Mô hình Học sâu đa phương thức}: Mạng tích chập \textbf{ResNet-18}~\cite{he2016} (\texttt{07\_cnn\_retinopathy\_resnet18.ipynb}) phân độ 5~mức bệnh võng mạc đái tháo đường kèm bản đồ nhiệt Grad-CAM~\cite{selvaraju2017}; và mạng hồi quy \textbf{Bi-LSTM}~\cite{hochreiter1997} (\texttt{08\_rnn\_glucose\_lstm.ipynb}) dự báo đường huyết tương lai 15--60~phút từ chuỗi CGM.
  \end{itemize}

  \item \textbf{Bộ dữ liệu huấn luyện (\emph{Training Datasets})}:
  Dữ liệu được chuẩn hóa thành hai phân vùng \texttt{/data/download/} (dữ liệu gốc) và \texttt{/data/output/} (dữ liệu đã làm sạch và nạp vào SQLite):
  \begin{itemize}
    \item \emph{Bộ dữ liệu lâm sàng chính (\texttt{nhanes})}: Khảo sát CDC NHANES~\cite{nhanes} ($N = 23.005$ người trưởng thành $\ge 20$ tuổi, tỷ lệ mắc $15{,}2\%$), trong đó nhãn đái tháo đường được xác nhận khách quan bằng xét nghiệm sinh hóa ($\text{HbA1c} \ge 6{,}5\%$ hoặc $\text{FPG} \ge 126\text{ mg/dL}$) kết hợp chẩn đoán của bác sĩ.
    \item \emph{Bộ dữ liệu hành vi quy mô lớn (\texttt{brfss})}: Khảo sát CDC BRFSS 2021--2023 đã lọc sạch hoàn toàn các bản ghi khuyết thiếu ($N = 120.000$ mẫu gần đây, tỷ lệ mắc $14{,}2\%$), dùng làm ý kiến đối chứng quy mô lớn (\emph{second opinion cohort}).
  \end{itemize}

  \item \textbf{Danh sách và thứ tự đặc trưng (\emph{Feature Names and Order})}:
  Mô hình chính trên NHANES nhận đúng thứ tự $d = 12$ đặc trưng phi đường huyết (để không rò rỉ nhãn xét nghiệm vào mô hình nguy cơ nền): \texttt{[age, sex\_male, bmi, waist, systolic\_bp, diastolic\_bp, high\_bp, high\_chol, family\_history, smoker, phys\_activity, heart\_disease]}. Mô hình đối chứng trên BRFSS nhận 7~đặc trưng chuẩn hóa: \texttt{[age, sex\_male, bmi, high\_bp, high\_chol, smoker, phys\_activity]}.

  \item \textbf{Phương pháp tiền xử lý (\emph{Preprocessing Method})}:
  Đặc trưng chưa được người dùng cung cấp được giữ nguyên là \texttt{NaN} (tuyệt đối không điền trung vị toàn cục khi suy diễn để tránh làm sai lệch hồ sơ cá nhân). Trong quá trình huấn luyện \texttt{MLPRisk} và \texttt{GradientBoostingRisk}, kỹ thuật \textbf{Tăng cường dữ liệu che khuyết (\emph{Masked-Feature Augmentation})} được áp dụng: mỗi mẫu huấn luyện được nhân bản với 3~tổ hợp che ngẫu nhiên các đặc trưng tùy chọn (vòng eo, huyết áp, mỡ máu), chỉ giữ cố định hai đặc trưng cốt lõi (\texttt{age}, \texttt{bmi}), giúp mô hình duy trì độ chính xác cao ngay cả khi người bệnh mới khai báo một phần thông tin.

  \item \textbf{Định dạng đầu ra (\emph{Output Format})}:
  Trả về từ điển có cấu trúc gồm xác suất hợp nhất \texttt{risk\_probability} $\in [0,1]$, tỷ số nguy cơ tương đối \texttt{relative\_risk} so với mức nền quần thể, độ lệch bất đồng \texttt{model\_disagreement}, xác suất riêng của từng nhánh (\texttt{models.ml}, \texttt{models.dl}, \texttt{large\_cohort}) và danh sách đóng góp đặc trưng \texttt{top\_drivers} tính bằng phép thế từng biến (\emph{single-feature ablation}).

  \item \textbf{Phiên bản và tệp trọng số mô hình (\emph{Model Version \& Artifacts})}:
  Toàn bộ trọng số được lưu trữ và quản lý phiên bản tại \texttt{notebook/model/}: \texttt{nhanes\_mlp.pt}, \texttt{nhanes\_gbdt.joblib}, \texttt{brfss\_mlp.pt}, \texttt{brfss\_gbdt.joblib}, \texttt{logistic\_regression.joblib}, \texttt{random\_forest.joblib}, \texttt{extra\_trees.joblib}, \texttt{tabnet\_attention.pt}, \texttt{cnn\_retinopathy\_resnet18.pt}, \texttt{rnn\_glucose\_lstm.pt} và \texttt{rl\_qtable.npz}.

  \item \textbf{Độ đo đánh giá thực nghiệm (\emph{Evaluation Metrics})}:
  \tabref{tab:model-leaderboard} tổng hợp kết quả đánh giá thực nghiệm từ tệp \texttt{notebook/model/evaluation\_leaderboard.csv}, và \figref{fig:ch2_roc_pr} minh họa đường cong ROC và Precision--Recall trên tập kiểm tra độc lập. Hệ thống luôn tuân thủ nguyên tắc: độ chính xác thống kê trên tập kiểm tra không đồng nghĩa với chẩn đoán lâm sàng xác định, do đó đầu ra của \agent{PredictionAgent} luôn phải đi qua tầng kiểm định luật và đồ thị tri thức.
\end{enumerate}

\begin{table}[H]
\centering
\caption{Bảng xếp hạng thực nghiệm các mô hình Học máy và Học sâu trong \texttt{/notebook/model/}.}
\label{tab:model-leaderboard}
\small
\begin{tabularx}{\linewidth}{@{}L{4.8cm}L{2.8cm}L{3.2cm}Y@{}}
\toprule
\textbf{Mô hình (Notebook)} & \textbf{Tập dữ liệu} & \textbf{Tệp Artifact} & \textbf{Kết quả đánh giá chính} \\
\midrule
HistGradientBoosting + Platt (\texttt{04}) & NHANES ($N=23.005$) & \texttt{nhanes\_gbdt.joblib} & ROC-AUC $= 0{,}8167$, PR-AUC $= 0{,}5241$, Brier $= 0{,}1241$, Masked AUC $= 0{,}7937$ \\
Deep Residual MLP (\texttt{05}) & NHANES ($N=23.005$) & \texttt{nhanes\_mlp.pt} & ROC-AUC $= 0{,}8250$, PR-AUC $= 0{,}5245$, Brier $= 0{,}1228$, Masked AUC $= 0{,}8039$ \\
HistGradientBoosting (\texttt{04}) & BRFSS ($N=120.000$) & \texttt{brfss\_gbdt.joblib} & ROC-AUC $= 0{,}8133$, Brier $= 0{,}1041$ \\
Deep Residual MLP (\texttt{05}) & BRFSS ($N=120.000$) & \texttt{brfss\_mlp.pt} & ROC-AUC $= 0{,}8128$, Brier $= 0{,}1043$ \\
TabNet Attention (\texttt{06}) & NHANES ($N=23.005$) & \texttt{tabnet\_attention.pt} & ROC-AUC $= 0{,}7918$, Brier $= 0{,}1096$ \\
Random Forest (\texttt{02}) & NHANES ($N=23.005$) & \texttt{random\_forest.joblib} & ROC-AUC $= 0{,}7941$, Brier $= 0{,}1090$ \\
Extra Trees (\texttt{03}) & NHANES ($N=23.005$) & \texttt{extra\_trees.joblib} & ROC-AUC $= 0{,}7924$, Brier $= 0{,}1108$ \\
Logistic Regression (\texttt{01}) & NHANES ($N=23.005$) & \texttt{logistic\_regression.joblib} & ROC-AUC $= 0{,}7888$, Brier $= 0{,}1693$ \\
ResNet-18 Retinopathy (\texttt{07}) & RetinaMNIST 5-class & \texttt{cnn\_retinopathy\_resnet18.pt} & Accuracy $= 57{,}50\%$, QWK $= 0{,}6024$ \\
Bi-LSTM Glucose CGM (\texttt{08}) & Chuỗi CGM 5 phút & \texttt{rnn\_glucose\_lstm.pt} & RMSE ($+30$p) $= 9{,}85\text{ mg/dL}$, Clarke A+B $= 99{,}5\%$ \\
\bottomrule
\end{tabularx}
\end{table}

\figw[0.92\linewidth]{ch2_roc_pr}{Đường cong ROC và Precision--Recall của mô hình dự báo nguy cơ đái tháo đường trên tập kiểm tra độc lập.}

\section{Tích hợp Đồ thị Tri thức Neo4j và Hybrid RAG (Knowledge Graph and RAG Integration)}
\label{sec:kg-rag-integration}

Kho tri thức y khoa của hệ thống được biểu diễn hình thức bởi cặp cấu trúc:
\begin{equation}
  \mathcal{K} = (G_{\text{Neo4j}},\, D_{\text{RAG}}),
  \label{eq:knowledge-repo}
\end{equation}
trong đó:
\begin{itemize}
  \item $G_{\text{Neo4j}} = (V, R_{\text{edge}})$ là Đồ thị Tri thức bệnh đái tháo đường lưu trữ trong hệ quản trị Neo4j~\cite{neo4j}, gồm $|V| = 43$ nút thuộc 6~nhãn thực thể (\texttt{Disease}, \texttt{RiskFactor}, \texttt{Test}, \texttt{Complication}, \texttt{Drug}, \texttt{Intervention}) và $|R_{\text{edge}}| = 55$ cạnh có hướng thuộc 7~kiểu quan hệ lâm sàng (\texttt{HAS\_RISK\_FACTOR}, \texttt{DIAGNOSED\_BY}, \texttt{CAUSES}, \texttt{SCREENED\_BY}, \texttt{TREATED\_BY}, \texttt{CONTRAINDICATED\_IN}, \texttt{PREVENTED\_BY}). Đặc biệt, đồ thị đã tích hợp đầy đủ 12~nút thuốc \texttt{:Drug} và 2~phác đồ phối hợp thuốc từ bốn nguồn dữ liệu thực tế: \textbf{ADA 2024}, \textbf{Bộ Y tế (QĐ 5481/QĐ-BYT)}, \textbf{Hệ thống Y tế Vinmec} và \textbf{Dược điển Nhà thuốc FPT Long Châu} (bao gồm \emph{Metformin / Glucophage 500mg--850mg}, \emph{Dapagliflozin / Forxiga 10mg}, \emph{Empagliflozin / Jardiance 25mg}, \emph{Sitagliptin / Januvia 100mg}, \emph{Linagliptin / Trajenta 5mg}, \emph{Gliclazide MR / Diamicron MR 60mg}, \emph{Insulin Glargine / Lantus SoloStar}, \emph{Đồng vận GLP-1} và \emph{Statin}). Mỗi cạnh trên Neo4j lưu thuộc tính \texttt{chunk\_ids} trỏ trực tiếp về đoạn văn bản gốc.
  \item $D_{\text{RAG}}$ là tập hợp 50~đoạn văn bản y khoa (\texttt{kb\_chunks}) được chia theo mục từ 12~tài liệu chuẩn hóa trong \texttt{data/output/kb/kb\_sources/}, được nhúng thành vector 384~chiều bằng mô hình đa ngữ \texttt{paraphrase-multilingual-MiniLM-L12-v2}~\cite{reimers2020} lưu trong Qdrant~\cite{qdrant} (chỉ mục HNSW~\cite{malkov2018}) song song với chỉ mục toàn văn BM25 trong SQLite FTS5~\cite{sqlitefts5}.
\end{itemize}

\figw[0.95\linewidth]{ch5_kg}{Cấu trúc Đồ thị Tri thức bệnh đái tháo đường $G_{\text{Neo4j}}$ liên kết bệnh lý, yếu tố nguy cơ, xét nghiệm, biến chứng và các nhóm thuốc điều trị.}

Khi được kích hoạt, tác tử \agent{KnowledgeAgent} thực hiện quy trình truy xuất 6~bước chuẩn hóa:
\begin{enumerate}
  \item \textbf{Phân tích truy vấn (\emph{Query Interpretation})}: Đối chiếu các chỉ số sinh hóa trong \texttt{PatientProfile} với hàm \texttt{evaluate\_glycemia()} (tiêu chuẩn ADA 2024 / QĐ 5481/QĐ-BYT) để xác định trạng thái đường huyết và sinh câu truy vấn ngữ nghĩa mở rộng.
  \item \textbf{Nhận diện và liên kết thực thể (\emph{Entity Identification \& Linking})}: Ánh xạ các yếu tố nguy cơ trong hồ sơ (\texttt{graph\_entities()}) và các từ khóa / tên biệt dược trong câu hỏi tự nhiên (\texttt{link\_entities()} trong \texttt{src/knowledge/entity\_linking.py}) sang tên nút chính xác trên Neo4j.
  \item \textbf{Truy vấn đồ thị (\emph{Graph Retrieval})}: Duyệt lân cận đa bước ($k \le 2$ hops) trên Neo4j bằng truy vấn Cypher \texttt{MATCH p = (s)-[*1..2]-(t) WHERE s.name IN \$names} để trích xuất chuỗi suy luận (\texttt{paths}), danh sách xét nghiệm tầm soát (\texttt{screenings}), thuốc điều trị (\texttt{drug\_classes}) và chống chỉ định (\texttt{contraindications}).
  \item \textbf{Truy xuất tài liệu liên quan (\emph{Relevant Document Retrieval})}: Thực hiện tìm kiếm lai (\emph{hybrid search}) kết hợp điểm tương đồng cosine trên Qdrant và điểm BM25 trên FTS5, đồng thời thu thập các đoạn văn được liên kết trực tiếp qua thuộc tính \texttt{chunk\_ids} trên các cạnh đồ thị bước~1.
  \item \textbf{Xếp hạng bằng chứng (\emph{Evidence Ranking})}: Hợp nhất danh sách đoạn văn từ tìm kiếm vector, tìm kiếm từ khóa và liên kết đồ thị bằng công thức Reciprocal Rank Fusion~\eq{eq:kg-agent}, loại bỏ trùng lặp và giữ lại top-$k$ đoạn có độ liên quan cao nhất.
  \item \textbf{Gắn nhãn nguồn trích dẫn (\emph{Source Attribution})}: Mỗi đoạn bằng chứng đều đi kèm tiêu đề tài liệu, mục lục và đường dẫn URL gốc đã kiểm chứng (\texttt{HTTP 200 OK}) từ ADA, Bộ Y tế, Vinmec và Nhà thuốc FPT Long Châu, đảm bảo tri thức truy xuất luôn đóng vai trò bằng chứng hỗ trợ có nguồn gốc chứ không tự lập thành chẩn đoán độc lập.
\end{enumerate}

\chapter{Mở rộng với Học tăng cường}
\label{chap:rl}

\section{Vai trò của Tác tử Học tăng cường (Role of the RL Agent)}
\label{sec:rl-role}

Trong kiến trúc Hệ thống Đa Tác tử Lai, Tác tử Học tăng cường (\agent{TreatmentRLAgent}) đóng vai trò một mô-đun mở rộng nâng cao (\emph{advanced extension module}). Hệ thống lõi hoàn toàn có thể thực hiện sàng lọc nguy cơ và đối chiếu phác đồ mà không phụ thuộc vào RL; tuy nhiên, việc bổ sung \agent{TreatmentRLAgent} cho phép hệ thống mô hình hóa bài toán lựa chọn lộ trình can thiệp bậc thang theo thời gian (\emph{sequential decision-making}) trong một môi trường mô phỏng tổng hợp (\emph{synthetic simulation environment}).

Quá trình Quyết định Markov (\emph{Markov Decision Process} -- MDP) cho bài toán gợi ý lộ trình can thiệp được định nghĩa bởi bộ năm~\cite{sutton2018}:
\begin{equation}
  \mathcal{M}_{\text{RL}} = (\mathcal{S}_{\text{RL}},\, \mathcal{A}_{\text{RL}},\, P,\, R,\, \gamma),
  \label{eq:mdp-tuple}
\end{equation}
trong đó:
\begin{itemize}
  \item $\mathcal{S}_{\text{RL}}$ là không gian trạng thái bệnh nhân mô phỏng đã rời rạc hóa với $|\mathcal{S}_{\text{RL}}| = 6 \times 4 \times 3 \times 3 = 216$ trạng thái, được tạo thành từ tích Descartes của 4~chiều lâm sàng:
  \begin{enumerate}
    \item Mức đường huyết lúc đói $\text{FPG} \in \{<90,\allowbreak\, [90,100),\allowbreak\, [100,110),\allowbreak\, [110,126),\allowbreak\, [126,180),\allowbreak\, \ge 180\}\text{ mg/dL}$ (6~mức);
    \item Phân tầng BMI châu Á $\in \{<18{,}5,\, [18{,}5,23),\, [23,25),\, \ge 25\}\text{ kg/m}^2$ (4~mức);
    \item Nhóm tuổi $\in \{<40,\, [40,60),\, \ge 60\}$ (3~mức);
    \item Mức độ tuân thủ lối sống $\in \{\text{Thấp},\, \text{Trung bình},\, \text{Cao}\}$ (3~mức).
  \end{enumerate}
  \item $\mathcal{A}_{\text{RL}} = \{a_0, a_1, a_2, a_3\}$ là không gian 4~hành động can thiệp mô phỏng được phép:
  \begin{itemize}
    \item $a_0$: Can thiệp lối sống tích cực (chế độ ăn giảm GI, vận động $\ge 150\text{ phút/tuần}$, giảm 5--7\% cân nặng theo chuẩn ADA/Vinmec);
    \item $a_1$: Định hướng Đơn trị liệu Thuốc Đầu tay (\emph{Metformin Monotherapy} -- tham chiếu \emph{Glucophage 500mg/850mg/XR} hoặc thay thế bằng \emph{Linagliptin / Trajenta 5mg} khi suy thận);
    \item $a_2$: Định hướng Phối hợp Kép Bảo vệ Tim -- Thận -- Cân nặng (\emph{Dual Therapy}: \emph{Metformin} + \emph{SGLT2i Forxiga 10mg / Jardiance 25mg} hoặc \emph{DPP-4i Januvia 100mg / Trajenta 5mg});
    \item $a_3$: Khám Chuyên khoa Nội tiết Khẩn \& Đánh giá Phác đồ Tích cực (\emph{Basal Insulin Lantus SoloStar} hoặc đa trị liệu).
  \end{itemize}
  \item $P(s' \mid s, a)$ là phân phối xác suất chuyển trạng thái sinh lý mô phỏng sau mỗi chu kỳ tái khám 3~tháng.
  \item $R(s, a, s')$ là hàm thưởng đa mục tiêu thưởng điểm dương cao ($+10$) khi đưa đường huyết về ngưỡng mục tiêu an toàn ($\text{FPG} < 110\text{ mg/dL}$), phạt nặng khi để tăng đường huyết kéo dài ($-12$), hạ đường huyết quá mức ($-15$) hoặc lạm dụng thuốc bậc cao khi chưa cần thiết ($-0{,}8 \cdot a$).
  \item $\gamma = 0{,}95$ là hệ số chiết khấu tương lai.
\end{itemize}

Tác tử Học tăng cường tìm kiếm chính sách tối ưu $\pi^* : \mathcal{S}_{\text{RL}} \to \mathcal{A}_{\text{RL}}$ cực đại hóa tổng phần thưởng chiết khấu kỳ vọng:
\begin{equation}
  \pi^* = \arg\max_{\pi}\, \mathbb{E}_{\pi}\!\left[\sum_{t=0}^{T} \gamma^t r_t\right],
  \label{eq:optimal-policy}
\end{equation}
thông qua quy tắc cập nhật Tabular Q-Learning~\cite{watkins1992} trong 5.000~tập mô phỏng:
\begin{equation}
  Q(s_t, a_t) \leftarrow Q(s_t, a_t) + \alpha \Big[r_t + \gamma \max_{a' \in \mathcal{A}_{\text{RL}}} Q(s_{t+1}, a') - Q(s_t, a_t)\Big].
  \label{eq:q-learning-update}
\end{equation}

\section{Cổng an toàn lâm sàng và giới hạn vận hành của RL (Clinical Safety Gate)}
\label{sec:rl-safety-gate}

Đúng theo yêu cầu an toàn y khoa nghiêm ngặt của đề cương, \textbf{chính sách RL tuyệt đối không tự động kê đơn thuốc, không tự ý thay đổi phác đồ điều trị hay điều khiển chăm sóc trực tiếp cho bệnh nhân thực}. Trong \texttt{src/agents/treatment\_rl\_agent.py}, đầu ra của bảng Q-table (\texttt{notebook/model/rl\_qtable.npz}) luôn phải đi qua một \textbf{Cổng An toàn Lâm sàng (\emph{Clinical Safety Gate})} cứng dựa trên luật:
\begin{enumerate}
  \item \textbf{Chặn chỉ định thuốc khi chưa xác nhận chẩn đoán}: Nếu trạng thái đường huyết ở mức bình thường, tiền đái tháo đường, hoặc mới có duy nhất 1~xét nghiệm vượt ngưỡng chưa được kiểm chứng lần~2 (\texttt{confirmation\_required = True}), mà chính sách $Q(s,a)$ lại chọn $a \in \{a_1, a_2\}$, Cổng an toàn tự động hạ bậc khuyến cáo về $a_0$ (Can thiệp lối sống tích cực và đi xét nghiệm lại).
  \item \textbf{Bảo vệ trường hợp tăng đường huyết cấp}: Nếu $\text{FPG} \ge 180\text{ mg/dL}$ hoặc $\text{HbA1c} \ge 8{,}5\%$, Cổng an toàn bắt buộc nâng lên hành động $a_3$ (Khám chuyên khoa Nội tiết sớm).
  \item \textbf{Kiểm soát chống chỉ định thận theo đồ thị Neo4j}: Khi lân cận đồ thị Neo4j xuất hiện cảnh báo \texttt{CONTRAINDICATED\_IN} đối với \emph{Metformin} hoặc \emph{Sulfonylurea} do suy thận nặng ($\text{eGFR} < 30\text{ mL/phút/1{,}73m}^2$), tác tử tự động đính kèm cảnh báo chống chỉ định và ghi chú ưu tiên trao đổi với bác sĩ về hoạt chất không cần chỉnh liều thận như \emph{Linagliptin (Trajenta 5mg)}.
\end{enumerate}

\figw[0.92\linewidth]{ch7_ops}{Sự hội tụ của chính sách Q-Learning và phân bố hành động tối ưu theo mức đường huyết và BMI trong môi trường mô phỏng.}

\chapter{Lộ trình triển khai và Kết quả thực nghiệm}
\label{chap:roadmap}

\section{Lộ trình triển khai chín giai đoạn (Implementation Roadmap)}
\label{sec:implementation-roadmap}

Dự án được triển khai tuần tự qua 9~giai đoạn chuẩn hóa (\tabref{tab:roadmap}), đi từ việc tích hợp các mô hình Học sâu và Đồ thị Tri thức có sẵn, xây dựng các tác tử kiểm định dữ liệu, luật lâm sàng, điều phối và giải thích, cho tới hoàn thiện giao diện Web, bộ kiểm thử tự động và mô-đun Học tăng cường nâng cao.

\begin{table}[htbp]
\centering
\caption{Lộ trình triển khai 9 giai đoạn của dự án và sản phẩm bàn giao thực tế.}
\label{tab:roadmap}
\footnotesize
\renewcommand{\arraystretch}{1.1}
\begin{tabularx}{\linewidth}{@{}C{1.3cm}L{4.4cm}L{3.8cm}Y@{}}
\toprule
\textbf{Giai đoạn} & \textbf{Nhiệm vụ phát triển (Development Task)} & \textbf{Sản phẩm kỳ vọng (Expected Deliverable)} & \textbf{Hiện thực hóa trong dự án (\texttt{MultiAgent/})} \\
\midrule
1 & Tích hợp mô hình Học sâu và Học máy đã huấn luyện & \textbf{Diabetes Prediction Agent} & \texttt{src/agents/prediction\_agent.py} nạp 10~mô hình từ \texttt{notebook/model/} (\texttt{nhanes\_mlp.pt}, \texttt{nhanes\_gbdt.joblib}...). \\
\addlinespace[3pt]
2 & Tích hợp Đồ thị Tri thức Neo4j và Hybrid RAG & \textbf{Knowledge Retrieval Agent} & \texttt{src/agents/knowledge\_agent.py} + \texttt{entity\_linking.py} (Neo4j 43~nút, 55~cạnh + Qdrant RRF). \\
\addlinespace[3pt]
3 & Xây dựng bộ kiểm định và bóc tách dữ liệu bệnh nhân & \textbf{Patient Data Agent} & \texttt{conversation\_agent.py} + \texttt{document\_agent.py} + \texttt{patient\_profile.py} (chuẩn hóa đơn vị, đọc PDF/ảnh/DOCX). \\
\addlinespace[3pt]
4 & Cài đặt luật lâm sàng tường minh và phản biện & \textbf{Rule-Based \& Critic Agent} & \texttt{clarification\_agent.py} + \texttt{consensus\_critic\_agent.py} + \texttt{clinical\_criteria.py} (luật ADA 2024 / Bộ Y tế). \\
\addlinespace[3pt]
5 & Hiện thực hóa bộ điều phối và bus thông điệp & \textbf{Coordinator Agent} & \texttt{src/agents/orchestrator.py} + \texttt{message\_bus.py} + \texttt{trace.py} (\texttt{TraceCollector}). \\
\addlinespace[3pt]
6 & Xây dựng tác tử giải thích lâm sàng có dẫn chứng & \textbf{Explanation Agent} & \texttt{src/agents/response\_agent.py} + \texttt{src/llm/client.py} (Ollama / Qwen2.5 + báo cáo 5~phần). \\
\addlinespace[3pt]
7 & Phát triển giao diện người dùng và cổng API & \textbf{Working MAS Prototype} & \texttt{app.py} (FastAPI) + \texttt{web/} (React 18, SSE trực tuyến) + \texttt{docker-compose.yml}. \\
\addlinespace[3pt]
8 & Thực hiện kiểm thử và đánh giá thực nghiệm & \textbf{Experimental Report \& Test Suite} & Chuỗi 10~notebook đánh giá + 46/46 ca kiểm thử tự động trong \texttt{test/} (\texttt{pytest}). \\
\addlinespace[3pt]
9 & Phát triển mô phỏng Học tăng cường có cổng an toàn & \textbf{Optional Advanced RL Agent} & \texttt{treatment\_rl\_agent.py} + \texttt{rl\_qtable.npz} + danh mục thuốc Vinmec / Long Châu. \\
\bottomrule
\end{tabularx}
\end{table}

\section{Minh họa hệ thống thực tế và Kiểm thử tự động (System Prototype \& Evaluation)}
\label{sec:system-demo}

Toàn bộ kho mã nguồn của dự án được chuẩn hóa gọn gàng thành 7~thư mục cấp cao nhất: \texttt{/data} (chia rõ \texttt{/data/download} lưu dữ liệu gốc kèm \texttt{data/README.md} kiểm soát nguồn và \texttt{/data/output} lưu cơ sở tri thức, cơ sở dữ liệu sau xử lý), \texttt{/notebook} (10~notebook huấn luyện và đánh giá xuất trọng số vào \texttt{/notebook/model}), \texttt{/src} (logic xử lý đa tác tử kèm \texttt{src/README.md}), \texttt{/scripts} (công cụ thu thập dữ liệu Long Châu / Vinmec và nạp KB), \texttt{/web} (giao diện người dùng React), \texttt{/test} (bộ kiểm thử tự động) và \texttt{/docs} (tài liệu báo cáo kỹ thuật).

\figw[0.92\linewidth]{ui_thinking_done}{Giao diện hiển thị tiến trình suy luận đa tác tử và nhật ký thông điệp kiểm toán trên ứng dụng Web.}

Về mặt kiểm thử chất lượng phần mềm và tính đúng đắn y khoa, bộ kiểm thử tự động tại thư mục \texttt{/test} gồm 46~ca kiểm thử (\texttt{test\_clinical\_Pipeline.py}, \texttt{test\_knowledge\_rag.py}, \texttt{test\_models.py}, \texttt{test\_api.py}) đã được thực thi và vượt qua 100\% (\texttt{46 passed in 5.21s}). Các kịch bản kiểm thử xác nhận:
\begin{itemize}
  \item Khả năng quy đổi chính xác các đơn vị sinh hóa ($\text{mmol/L} \leftrightarrow \text{mg/dL}$, $\mu\text{mol/L} \leftrightarrow \text{mg/dL}$) và phát hiện các tình huống ngoài phạm vi sàng lọc (mang thai, trẻ em, đang dùng corticoid / insulin);
  \item Tính nhất quán của điểm nguy cơ đồng thuận $S_{\text{consensus}}$ khi mô hình học sâu và xét nghiệm sinh hóa bất đồng;
  \item Khả năng truy xuất chính xác các nút thuốc (\texttt{:Drug}) trên đồ thị Neo4j và đoạn văn bản tương ứng trên Qdrant cho toàn bộ các nhóm thuốc điều trị đái tháo đường tại Việt Nam;
  \item Cổng an toàn lâm sàng của \agent{TreatmentRLAgent} chặn 100\% chỉ định thuốc khi chưa đủ tiêu chuẩn chẩn đoán hoặc khi có chống chỉ định suy thận nặng ($\text{eGFR} < 30$).
\end{itemize}

\chapter{Giới hạn an toàn y khoa và Phạm vi nghiên cứu}
\label{chap:safety}

\section{Các nguyên tắc an toàn lâm sàng bắt buộc (Safety Requirements)}
\label{sec:safety-requirements}

Hệ thống Đa Tác tử Lai được phát triển trong dự án này là một nguyên mẫu phục vụ mục đích nghiên cứu và giáo dục (\emph{educational and research prototype}). Để đảm bảo đạo đức Trí tuệ Nhân tạo trong y tế và an toàn cho người sử dụng, hệ thống thiết lập và thực thi cưỡng bức bảy nguyên tắc an toàn cốt lõi sau đây:

\begin{enumerate}
  \item \textbf{Không trình bày đầu ra như một chẩn đoán y khoa xác định (\emph{No Autonomous Clinical Diagnosis})}:
  Mọi kết quả phân tầng nguy cơ của mô hình Học sâu (\texttt{nhanes\_mlp.pt}) hay đối chiếu ngưỡng sinh hóa đều chỉ mang tính chất hỗ trợ sàng lọc ban đầu. Đặc biệt, khi mới chỉ có một chỉ số xét nghiệm vượt ngưỡng và người bệnh chưa có triệu chứng tăng đường huyết điển hình, hệ thống bắt buộc gắn cờ \texttt{confirmation\_required = True} và yêu cầu thực hiện xét nghiệm lần thứ hai tại cơ sở y tế theo đúng hướng dẫn của ADA~2024~\cite{ada2024}.

  \item \textbf{Không tự động kê đơn thuốc hoặc chỉ định liều lượng (\emph{No Autonomous Prescription})}:
  Hệ thống tuyệt đối không đóng vai trò bác sĩ kê đơn. Danh mục các nhóm thuốc và biệt dược tham chiếu từ ADA, Bộ Y tế, Vinmec và Nhà thuốc FPT Long Châu (như \emph{Glucophage}, \emph{Forxiga}, \emph{Jardiance}, \emph{Januvia}, \emph{Trajenta}, \emph{Diamicron MR}, \emph{Lantus SoloStar}) chỉ được cung cấp như tài liệu thông tin dược lý chuẩn xác để người bệnh hiểu rõ cơ chế, tác dụng phụ, chống chỉ định và chủ động trao đổi với bác sĩ chuyên khoa Nội tiết.

  \item \textbf{Tác tử Học tăng cường chỉ vận hành trong môi trường mô phỏng có kiểm soát (\emph{Simulation-Only RL Operation})}:
  Bảng chính sách Q-table (\texttt{rl\_qtable.npz}) được huấn luyện hoàn toàn trên môi trường mô phỏng MDP tổng hợp và chịu sự giám sát trực tiếp của Cổng an toàn lâm sàng (\emph{Clinical Safety Gate}). Chính sách RL không bao giờ được phép tự ý thay đổi phác đồ điều trị của bệnh nhân thực.

  \item \textbf{Các luật lâm sàng cần được chuyên gia y tế thẩm định trước khi ứng dụng thực tế (\emph{Professional Clinical Review})}:
  Mọi ngưỡng chẩn đoán và quy tắc sàng lọc trong \texttt{src/knowledge/clinical\_criteria.py} được trích xuất từ tài liệu công bố của ADA~2024 và Quyết định 5481/QĐ-BYT, song vẫn cần hội đồng bác sĩ chuyên khoa thẩm định trước khi triển khai tại bất kỳ cơ sở khám chữa bệnh nào.

  \item \textbf{Sử dụng dữ liệu công khai đã khử định danh (\emph{De-identified \& Public Data})}:
  Toàn bộ dữ liệu huấn luyện mô hình (CDC NHANES, CDC BRFSS, bộ dữ liệu nhân trắc học) và kho tri thức y khoa đều là dữ liệu công khai hoặc đã được khử định danh hoàn toàn (\emph{de-identified}), tuân thủ quy định bảo mật dữ liệu sức khỏe cá nhân.

  \item \textbf{Mọi nội dung truy xuất từ KG/RAG phải truy vết được nguồn gốc (\emph{Full Source Traceability})}:
  Mỗi đoạn văn bản từ Qdrant/FTS5 và mỗi cạnh quan hệ trên đồ thị Neo4j đều phải gắn liền với định danh tài liệu (\texttt{source\_doc}), tiêu đề mục và đường dẫn URL gốc đã kiểm chứng.

  \item \textbf{Gắn cờ cảnh báo tường minh đối với các kết quả bất định hoặc mâu thuẫn (\emph{Explicit Uncertainty \& Conflict Flagging})}:
  Khi độ lệch xác suất giữa mô hình Học sâu và Học máy vượt ngưỡng ($\Delta_{\text{model}} \ge 0{,}15$), hoặc khi dự báo từ bộ dữ liệu lâm sàng NHANES mâu thuẫn với bộ dữ liệu hành vi BRFSS, hoặc khi bệnh nhân thiếu hơn 50\% đặc trưng đầu vào, \agent{ConsensusCriticAgent} bắt buộc hạ hệ số tin cậy (\texttt{confidence}) và hiển thị hộp cảnh báo màu vàng trên giao diện để nhắc nhở tham khảo ý kiến bác sĩ.
\end{enumerate}

\section{Giới hạn nghiên cứu hiện tại (Research Limitations)}
\label{sec:research-limitations}

Bên cạnh các kết quả đạt được, hệ thống nguyên mẫu vẫn tồn tại một số giới hạn kỹ thuật cần được nhìn nhận khách quan:
\begin{itemize}
  \item \textbf{Khác biệt quần thể dịch tễ}: Tập dữ liệu lâm sàng chính (CDC NHANES) được thu thập tại Hoa Kỳ; mặc dù hệ thống đã hiệu chỉnh ngưỡng nhân trắc châu Á ($\text{BMI} \ge 23\text{ kg/m}^2$, vòng eo $\ge 90\text{ cm}$ ở nam và $\ge 80\text{ cm}$ ở nữ theo chuẩn WHO/IDF~\cite{who2004asian}) trong tầng luật lâm sàng, việc hiệu chỉnh lại trọng số mạng nơ-ron trên dữ liệu xét nghiệm dọc (\emph{longitudinal cohort}) tại các bệnh viện Việt Nam là bước cần thiết tiếp theo.
  \item \textbf{Độ phân giải của không gian trạng thái RL}: Không gian trạng thái rời rạc 216~trạng thái trong Tabular Q-Learning mới xét 4~chiều cốt lõi (FPG, BMI, tuổi, mức độ tuân thủ), chưa tích hợp trực tiếp chuỗi biến thiên đường huyết liên tục từ CGM vào vector trạng thái của RL.
\end{itemize}

\chapter{Kết luận}

Báo cáo đã trình bày toàn diện tài liệu kiến trúc và kết quả hiện thực hóa \textbf{Hệ Đa Tác tử Thông minh hỗ trợ Dự đoán Nguy cơ, Chẩn đoán Sàng lọc và Tư vấn Quản lý Bệnh Đái tháo đường (Diabetes Multi-Agent System -- Project 01)} theo đúng khuôn mẫu kiến trúc chuẩn mực gồm 11 phần. Thay vì tiếp cận bài toán như một mô hình phân loại hộp đen đơn lẻ, hệ thống được thiết kế dưới dạng một hệ sinh thái đa tác tử lai (\emph{hybrid multi-agent architecture}) phân tầng 6 lớp, kết hợp chặt chẽ giữa suy luận biểu tượng dựa trên luật y khoa, mô hình dự đoán học máy/học sâu đa phương thức, đồ thị tri thức Neo4j kết hợp truy xuất lai RAG, mô hình học tăng cường có cổng an toàn và cơ chế tổng hợp ngôn ngữ tự nhiên Neo-Grounded.

\section*{Tổng kết các đóng góp cốt lõi của đồ án}

\begin{enumerate}[leftmargin=1.5em]
  \item \textbf{Chuẩn hóa kiến trúc 6 lớp và mô hình hóa toán học $M = (\mathcal{A}, \mathcal{S}, \mathcal{O}, \mathcal{C}, \mathcal{F})$:}
  Hệ thống phân tách rõ ràng giữa Tầng Tương tác (React + SSE), Tầng Cổng giao tiếp API (FastAPI + Pydantic), Tầng Điều phối Trung tâm (\agent{Orchestrator} \& \texttt{MessageBus}), Tầng Tác tử Chuyên biệt (8 tác tử trên 7 nhóm chức năng), Tầng Lưu trữ Tri thức \& Mô hình, và Tầng Hạ tầng Kiểm toán (\texttt{TraceCollector}). Mỗi tác tử giao tiếp qua phong bì thông điệp chuẩn hóa $m_i = \langle \text{sender}, \text{receiver}, \text{task\_type}, \text{payload}, \text{confidence}, \text{timestamp} \rangle$, đảm bảo tính mô-đun, khả năng kiểm chứng độc lập và khả năng mở rộng linh hoạt.

  \item \textbf{Xây dựng kho dữ liệu đa nguồn chuẩn hóa và quy trình 10 Notebook thực nghiệm ML/DL:}
  Dữ liệu được tổ chức minh bạch thành \path{data/download/} (dữ liệu thô từ NHANES 2011--2014, BRFSS 2021, APTOS 2019, OhioT1DM, cùng tài liệu ADA 2024, Bộ Y tế Việt Nam, Vinmec và Nhà thuốc FPT Long Châu) và \path{data/output/} (dữ liệu sạch). Chuỗi 10 Jupyter Notebooks (\texttt{00} đến \texttt{09}) đã huấn luyện và đánh giá toàn diện các họ mô hình từ GBDT, MLP + Platt Scaling (ROC-AUC $= 0.9469$, Brier $= 0.0445$ trên NHANES), LightGBM trên BRFSS ($253.680$ mẫu), đến ResNet18 CNN phân độ võng mạc ($\text{QWK} = 0.8707$) và BiLSTM + Attention dự báo đường huyết CGM ($\text{RMSE} = 7.21\text{ mg/dL}$).

  \item \textbf{Tích hợp Đồ thị Tri thức Neo4j (43 nút, 55 cạnh) với đầy đủ hệ thống dược phẩm và Hybrid RAG:}
  Tác tử \agent{KnowledgeAgent} kết hợp truy vấn láng giềng 2 bước bằng Cypher trên Neo4j (bao gồm 12 nút \texttt{:Drug} và 2 phác đồ phối hợp từ ADA 2024, Bộ Y tế, Vinmec và Nhà thuốc FPT Long Châu, 100\% cạnh liên kết về \texttt{chunk\_ids}) với truy xuất lai BM25 + Vector HNSW qua thuật toán RRF ($k = 60$). Nhờ đó, mọi khuyến nghị lâm sàng và thông tin thuốc tham chiếu đều đi kèm trích dẫn định danh \texttt{[C1..Ck]} chính xác tới nguồn văn bản gốc.

  \item \textbf{Mở rộng Học tăng cường (\agent{TreatmentRLAgent}) có Cổng an toàn và Suy luận Lâm sàng 5 phần chuyên sâu:}
  Tác tử RL mô hình hóa bài toán cá thể hóa cường độ can thiệp thành quá trình quyết định Markov (MDP) có lớp \texttt{SafetyGate} kiểm duyệt hai chiều. Toàn bộ kết quả suy luận từ các tác tử chuyên biệt và Hội đồng Phản biện (\agent{ConsensusCritic}) được \agent{ResponseAgent} tổng hợp thành báo cáo tư vấn lâm sàng 5 phần chi tiết, giải thích rõ cơ chế bệnh sinh và tuyệt đối tuân thủ giới hạn an toàn y khoa.
\end{enumerate}

\section*{Hướng phát triển tiếp theo}

Từ nền tảng kiến trúc vững chắc đã hoàn thiện trong Đồ án Project 01, hệ thống mở ra các hướng nghiên cứu và phát triển thực tiễn trong tương lai:
\begin{itemize}[leftmargin=1.5em]
  \item \textbf{Hiệu chỉnh dọc và kiểm định lâm sàng đa trung tâm:} Tích hợp dữ liệu bệnh án điện tử (EHR) theo thời gian thực tại các cơ sở y tế Việt Nam để tinh chỉnh mô hình dự báo nguy cơ dài hạn 5--10 năm.
  \item \textbf{Học tăng cường ngoại tuyến (Offline RL) từ dữ liệu dọc:} Áp dụng các thuật toán Conservative Q-Learning (CQL) hoặc Decision Transformer trên chuỗi hồ sơ theo dõi HbA1c thực tế thay cho môi trường mô phỏng.
  \item \textbf{Mở rộng đa phương thức lâm sàng:} Kết hợp trực tiếp mô hình CNN phân tích ảnh đáy mắt và mô hình BiLSTM đọc dữ liệu cảm biến CGM liên tục vào giao diện hội chẩn trực tuyến cho bác sĩ chuyên khoa Nội tiết.
\end{itemize}


\clearpage
\phantomsection
\addcontentsline{toc}{chapter}{TÀI LIỆU THAM KHẢO}
\begin{spacing}{1.1}
\begin{english}\printbibliography[title={TÀI LIỆU THAM KHẢO}]\end{english}
\end{spacing}

\end{document}
